\documentclass[12pt,a4paper]{article}
\usepackage[T2A]{fontenc}
\usepackage[utf8]{inputenc}
\usepackage[english,russian]{babel}
\usepackage{amsmath,amssymb,amsfonts}
\usepackage{geometry}
\usepackage{enumitem}
\usepackage{booktabs}
\usepackage{longtable}
\usepackage{array}
\usepackage[hidelinks]{hyperref}

\geometry{margin=2.5cm}
\setlength{\parskip}{0.6ex}
\setlength{\parindent}{0pt}
\emergencystretch=2em
\widowpenalty=1000
\clubpenalty=1000

\newcommand{\id}[1]{\texttt{#1}}
\newcommand{\mech}[1]{\paragraph{Механизм.} #1}
\newcommand{\why}[1]{\paragraph{Почему так.} #1}
\newcommand{\alt}[1]{\paragraph{Альтернативы.} #1}

\title{Логика модели: обоснование решений и альтернативы\\[0.4em]
\large Микросимуляционная модель скрининга онкологических заболеваний (MedicalModel2024)}
\author{Документация модели}
\date{\today}

\begin{document}
\maketitle
\tableofcontents
\newpage

\section{Назначение, рамки и как читать этот документ}

Этот документ описывает, \emph{что делает модель и почему она устроена именно так}. Он сознательно не
является ни руководством по коду, ни обзором переноса. Почти для каждого элемента модели существует более
одного защитимого способа реализации, поэтому полезная информация --- это не только сделанный выбор, но и
отвергнутые альтернативы вместе с причинами отказа: именно эти причины определяют, где модель сильна и где
другие данные или другой вопрос потребовали бы иной конструкции.

Каждый элемент модели описан по одному и тому же шаблону из трёх частей:

\begin{enumerate}[nosep]
  \item \textbf{Механизм} --- правило, которое модель фактически применяет, изложенное так, чтобы его можно
        было воспроизвести по одному этому описанию.
  \item \textbf{Почему так} --- что даёт выбранное решение: идентифицируемость по доступным данным,
        эпидемиологическую правдоподобность, интерпретируемость или вычислительную стоимость.
  \item \textbf{Альтернативы} --- другие защитимые конструкции, причины, по которым каждая не использована,
        и что должно измениться, чтобы её применить.
\end{enumerate}

Два смежных документа покрывают остальные ракурсы и здесь не повторяются:

\begin{itemize}[nosep]
  \item \id{model\_for\_dummies.pdf} --- изложение модели и веб-интерфейса на общедоступном языке.
  \item \id{python\_port\_analysis.pdf} --- технический и статистический разбор: переменные состояния,
        статистические исправления, производительность и методологические улучшения.
\end{itemize}

\textbf{Эталонная конфигурация.} Все числа, приводимые ниже как иллюстрации, взяты из одной конфигурации:
наборы данных GLOBOCAN по колоректальному раку (агрегированный файл США и файл стадий), $N = 20\,000$
агентов, зерно 2022, 30 моделируемых лет, скрининг в возрасте 50--60 лет с периодичностью два года и
участием 50\%. Такой прогон даёт 1\,142 случая рака, 200 пациентов, излеченных клинически, 5 излеченных
после обнаружения при скрининге, 12\,334 выживших к концу прогона и суммарные затраты 13\,509 единиц.
Числа из других прогонов напрямую не сопоставимы, поскольку модель стохастична.

\textbf{Об честности изложения.} Там, где элемент присутствует в коде, но не влияет на результаты, он прямо
перечислен в разделе~\ref{sec:inactive}, а не замалчивается. Там, где модель \emph{не может} определить
величину по доступным данным, это прямо сказано в том месте, где это важно.

\section{Класс модели, время и рамки популяции}

\subsection{Почему микросимуляция на агентах}

\mech{Модель создаёт $N$ независимых агентов. Каждому агенту выдаётся одна полная история рака ---
возникает ли первая злокачественная клетка и когда, как быстро растёт опухоль, какая стадия достигнута к
моменту клинической диагностики, излечивает ли лечение и в каком возрасте человек умирает, --- после чего
агент прослеживается год за годом. Ничего не агрегируется в доли; популяционные показатели получаются
подсчётом агентов постфактум.}

\why{Такая конструкция навязана четырьмя свойствами самого вопроса.

\emph{(i) Ответ --- это разность хвостов, а не среднее.} Скрининг оправдан небольшим числом жизней,
выведенных из колонки смертей от рака, и числом лет жизни, которые эти люди получают. И то и другое ---
свойства индивидуальных контрфактических траекторий, а именно их модель средних представлять не может:
один и тот же человек должен наблюдаться и с вмешательством, и без него. Поэтому модель хранит
\emph{оба} возраста смерти для каждого пациента (без лечения и с учётом скрининга), и кривая
дополнительно выживших строится из внутрииндивидуальных разностей, а не из двух независимых популяций.
Неоднородность и хвост неразделимы: выигрывают не средние люди.

\emph{(ii) Скрининг зависит от истории.} Кто будет обследован в 58 лет, зависит от участия в 56, от того,
не найден ли уже рак, и от возрастного окна. В модели состояний это заставляет нести всю историю
скрининга в пространстве состояний. В микросимуляции это несколько булевых признаков на агента.

\emph{(iii) Неоднородность здесь структурна, а не декоративна.} Индивидуальная скорость роста
($B_{\text{ind}}$), агрессивность по стадиям, зависящая от возраста эффективность излечения и
эффективность лечения по стадиям действуют мультипликативно на исход. Достоверно представить любую из них
в когортной модели --- значит умножить число состояний на число уровней; их совместный учёт взрывоопасен
комбинаторно.

\emph{(iv) Модель должна воспроизводить распределение по стадиям, а не среднюю стадию.} Стадия при
диагностике есть следствие того, как долго растёт опухоль и как быстро, то есть по своей природе ---
событие пересечения порога конкретным агентом (раздел~\ref{sec:stage}).}

\alt{\begin{itemize}[nosep]
  \item \textbf{Детерминированная когортная марковская модель / многосостоянийная таблица дожития.}
        Классическая альтернатива, применяемая во многих анализах скрининга, поскольку она быстра,
        детерминирована и проста для изложения. Отвергнута, потому что не может нести индивидуальный
        контрфактический сценарий, которым и определяется выход «число спасённых жизней», и потому что
        пространство состояний растёт мультипликативно с неоднородностью.
  \item \textbf{Возрастно-структурированное уравнение в частных производных (Маккендрик --- фон Фёрстер).}
        Аналитически разрешимо и быстро для вопросов о средних. Отвергнуто по той же причине плюс по
        практической: скрининг делает скорости переходов зависящими от времени и истории, что превращает
        уравнение в большую систему.
  \item \textbf{Уравнение восстановления / интегральная проекционная модель.} Очень эффективно для
        заболеваемости и сдвига стадий. Отвергнуто, потому что смертность от конкретной причины и
        внутрииндивидуальный обмен смерти от рака на смерть от других причин в этой форме выражаются
        неудобно.
  \item \textbf{Полумарковский процесс с распределениями времён пребывания.} Хороший компромисс и реальная
        опция. Здесь отвергнут, потому что механизм пересечения стадий уже является полумарковским
        процессом с явной моделью роста; добавление второй латентной структуры прогрессии дало бы двойной
        учёт.
  \item \textbf{Компартментные ОДУ / системная динамика.} Отвергнуто: нет индивидуальной изменчивости, нет
        истории скрининга, а вопрос «кто выигрывает» становится безответным.
  \item \textbf{Дерево решений или байесовская сеть.} Уместны для однократного решения, а не для
        пожизненного вмешательства с конкурирующими рисками.
\end{itemize}}

\subsection{Почему шаг времени равен одному году}

\mech{Время --- это целочисленный счётчик лет \id{current\_date}. Возраст целочисленный
($\text{возраст} = \id{current\_date} - \id{date\_birth}$). На каждом шаге модель (а) удаляет агентов, чей
возраст смерти равен текущему возрасту, (б) проводит раунд скрининга, если год попадает в сетку скрининга,
и (в) обновляет статистику. Таким образом, все события внутри года разрешаются на границах года.}

\why{Шаг времени выбран под гранулярность данных калибровки и вмешательства, а не для удобства. Агрегированные
данные --- это годовые таблицы по возрастным группам; оцениваемая программа задана на годовой сетке
(«приглашение в 50--60 лет раз в два года»); а вовлечённые процессы естественной истории медленны по
сравнению с годом, с одним исключением, о котором ниже. Целочисленный возраст делает три внутренних
механизма точными, а не приближёнными: сравнение возрастов смерти становится равенством целых, стадия при
обнаружении при скрининге --- поиском сохранённого возраста пересечения относительно текущего возраста, а
экспортируемая таблица агентов --- читаемой напрямую.}

\why{Цена названа прямо: внутри года события одновременны и неупорядочены, ни одно событие не может
произойти дважды за год, а процесс, более быстрый чем год --- терминальная фаза после последнего
пересечения стадии, --- дискретизируется. Последнее здесь не сноска: подобранный риск смерти от рака даёт
среднюю терминальную выживаемость около трёх лет (раздел~\ref{sec:calib}), поэтому годовая сетка адекватна,
но не с запасом.}

\alt{\begin{itemize}[nosep]
  \item \textbf{Непрерывное время с экспоненциальными часами (схема «следующего события» в духе
        Джиллеспи).} Точная динамика внутри года и точные моменты скрининга; требует экспоненциального
        предположения для каждого перехода и существенно большего объёма реализации. Это наиболее
        защитимое улучшение, если предметом изучения когда-нибудь станет именно темп терминальной фазы.
  \item \textbf{Месячный или квартальный шаг.} Точнее, но умножает стоимость на 12 или 4, тогда как
        калибровочные данные не позволяют задать внутригодовые цели, поэтому дополнительное разрешение не
        будет контролироваться данными.
  \item \textbf{Событийное моделирование с очередью приоритетов.} Естественно для непрерывного времени;
        избыточная механика для годовой сетки.
  \item \textbf{Гибрид: годовая сетка скрининга при непрерывной естественной истории.} Разумный
        компромисс, но он усложняет воспроизводимость и экспортируемую таблицу агентов без выигрыша при
        текущем уровне данных.
\end{itemize}}

\subsection{Почему когорта замкнута, а горизонт равен 30 годам}

\mech{Рождений и миграции нет: популяция только взрослеет и умирает. Начальная возрастная структура
берётся из наблюдаемой популяции, поэтому когорта уже содержит все возрасты при $t=0$. Горизонт задаётся
параметром \id{years\_to\_simulate} (по умолчанию 30).}

\why{Замкнутая когорта изолирует оцениваемый эффект. Если бы моделировались рождения и миграция,
наблюдаемая разность смертности между сценариями со скринингом и без него смешала бы эффект скрининга с
демографическим дрейфом, а модель нуждалась бы в предположениях о рождаемости и миграции, которые ей нечем
проверить. Поскольку начальное распределение возрастов поперечное --- люди входят в модель в своём текущем
возрасте, --- когорта реалистична уже при $t=0$, прогрев не нужен, и первый год выпуска данных пригоден к
использованию; это важно для окна скрининга, открывающегося в 50 лет. Тридцати лет хватает, чтобы дважды
покрыть всё окно скрининга и накопить разность смертности, сохраняя прогон интерактивным: эталонный прогон
на 20\,000 агентов и 30 лет завершается за пару секунд.}

\alt{\begin{itemize}[nosep]
  \item \textbf{Когорта рождения, прослеженная с возраста 0.} Наиболее чистая конструкция для пожизненных
        исходов и взаимодействий с детскими экспозициями; требует прогрева длиной в столетие и значительно
        больше агентов для той же статистической мощности в возрастах 50--80, где и принимаются решения.
  \item \textbf{Открытая популяция с демографическим прогнозом.} Реалистичнее для долгосрочных прогнозов;
        вводит предположения о рождаемости и миграции, которые модель не может подкрепить данными, и
        загрязняет контраст по скринингу.
  \item \textbf{Стационарная популяция.} Привлекательна, если интересны только долгосрочные показатели;
        добавляет предположение о демографическом равновесии, не помогая вопросу о скрининге.
  \item \textbf{Гибридный вход в середине жизни.} Вход в когорту в 40--50 лет ещё сократил бы прогоны, но
        отбросил бы взаимодействие возраста начала скрининга с историей до скрининга.
\end{itemize}}

\section{Демография: начальная возрастная структура и конкурирующая смертность}

\subsection{Начальная возрастная структура}

\mech{Каждому агенту присваивается начальный возраст выборкой из наблюдаемой возрастной структуры
популяции: колонка \id{population} агрегированного файла, по одному значению на год возраста, нормируется
и используется как дискретное распределение; возрасты разыгрываются методом обратной функции распределения
(равномерная величина локализуется в накопленном массиве). Затем год рождения агента задаётся как минус
этот возраст, так что абсолютный возраст --- это просто «время с момента рождения».}

\why{Поперечная возрастная структура --- правильное начальное условие для замкнутой когорты: она
фиксирует и численность группы риска, и возрастной профиль этого риска в момент старта симуляции.
Использование файла как распределения вероятностей, а не как абсолютных численностей, означает, что
единицы измерения файла сокращаются, и модель не зависит от того, выражена колонка в человеках, тысячах
или процентах. Годы возраста взяты поштучно, потому что данные уже есть в таком разрешении; модель не
навязывает собственного группирования, поэтому любая потеря детализации --- выбор данных, а не модели.}

\alt{\begin{itemize}[nosep]
  \item \textbf{Равномерное или одновершинное распределение возрастов.} Отвергнуто: возрастное
        распределение --- главный драйвер и заболеваемости, и смертности, поэтому нереалистичное
        распределение испортило бы все выходные показатели.
  \item \textbf{Возрастные группы ступенчатой функцией.} Дешевле и ближе к тому, как часто публикуются
        исходные данные; теряется разрешение внутри группы, что существенно рядом с окном скрининга.
  \item \textbf{Подогнанный параметрический профиль возрастов (Гомпертц, Вейбулл, сплайн).} Сглаживает
        выборочный шум и умеет экстраполировать, но навязывает форму, которой данные не требуют, и любое
        расхождение распространяется на все показатели.
  \item \textbf{Эмпирические микроданные} (выборка переписи или обследования с индивидуальными возрастами).
        Самая точная опция; требует данных, которых в комплекте модели нет, тогда как агрегированный путь и
        так воспроизводит маргинальное распределение.
  \item \textbf{Старт всех агентов с одного возраста} (например, все по 50). Отвергнуто: разрушает
        возрастной состав, удаляет самых молодых и самых старых и превращает окно скрининга в вырожденную
        границу когорты.
\end{itemize}}

\subsection{Смертность не от рака как конкурирующий риск}

\mech{Каждый агент также получает \emph{естественный возраст смерти}, разыгрываемый тем же способом из
распределения, построенного по наблюдаемым годовым числам смертей от всех причин (\id{deaths all},
делённым на общую численность популяции и нормированным). Агенты, чей разыгранный возраст смерти не
превышает их начального возраста, разыгрываются заново один раз; если и вторая попытка не проходит, возраст
смерти полагается равным начальному возрасту плюс один, то есть агент умирает в первый же моделируемый год.
Далее естественный возраст смерти фиксирован, и всё, что делает модель, сравнивается с ним.}

\why{Представление смертности не от рака единым фиксированным индивидуальным «возрастом смерти»,
разыгранным один раз, даёт три вещи.

\emph{(i) Конкурирующий риск превращается в арифметическое сравнение.} Имея на агента один естественный
возраст смерти и один возраст смерти от рака, всё, что зависит от их взаимодействия --- что убило человека
раньше, излечен ли он, спас ли его скрининг, сколько лет жизни он получил, --- есть сравнение двух чисел.
Именно это делает парный учёт «скрининг против отсутствия скрининга» точным, а не приближённым: контрольную
популяцию не нужно перезапускать, и в контраст не добавляется выборочный шум.

\emph{(ii) Это эквивалентно модели выживаемости, но без сопутствующего учёта.} Разыгрывание возраста смерти
из таблицы дожития маргинально эквивалентно ежегодному удалению доли выживших, но заменяет текущую
переменную состояния и итеративное обновление одним сохранённым целым числом.

\emph{(iii) Условный пересев --- правильная поправка.} Поскольку агенты входят в модель в своём текущем
возрасте, их возраст смерти должен быть условным по факту того, что они уже живы; повторный розыгрыш
меньшинства, у которого первая попытка оказалась ниже возраста входа, реализует именно это усечение, а
резервный вариант (возраст $+1$) гарантирует корректную жизнь длиной в год вместо нулевой или
отрицательной.}

\why{\emph{Следствие, которое нужно сформулировать прямо.} Агент удаляется из популяции только по
достижении естественного возраста смерти. Зафиксированная смерть от рака агента \emph{не} удаляет: смерти от
рака считаются событием конкретной причины в слое статистики. Поэтому возрастная структура точно следует
таблице дожития по всем причинам, а знаменатель группы риска для показателей по причинам всё ещё включает
агентов, у которых смерть от рака уже зафиксирована. Это стандартный приём оценивания риска конкретной
причины со знаменателем общей популяции, и он корректен при обычном предположении о независимости причин
--- что удаление одной причины не меняет силу другой. Он был бы неверен, если бы считалось, что смерть от
рака меняет последующий риск смерти от других причин; это одно из предположений, которое должна проверять
валидация.}

\alt{\begin{itemize}[nosep]
  \item \textbf{Ежегодный розыгрыш Бернулли с признаком выживания.} Маргинально эквивалентен для одной
        причины, но делает возраст смерти выходом, а не входом, что резко усложняет парный контрфактический
        сценарий, и добавляет переменную состояния, которую нужно согласованно поддерживать.
  \item \textbf{Параметрический закон смертности} (Гомпертца --- Мейкхэма, Хелигмана --- Полларда,
        Ли --- Картера). Гладкий, компактный, умеет экстраполировать; здесь не нужен, поскольку наблюдается
        весь диапазон возрастов (0--110), так что подогнанный закон мог бы лишь добавить расхождение.
  \item \textbf{Таблицы дожития с удалённой причиной (прочие причины).} Концептуально правильная
        конструкция для модели рака: удалить смерти от рака из фона, чтобы две причины не учитывались
        дважды. Не используется, потому что в комплекте есть только смерти от всех причин, а нужен столбец
        смертей по возрастным группам для конкретной причины. Вспомогательная функция перевода годовых
        рисков в распределение возраста смерти (\id{risks\_to\_distribution},
        $P(T=a)=q(a)\prod_{j<a}(1-q(j))$) присутствует в коде, но никогда не вызывается: модель использует
        наблюдаемые смерти напрямую.
  \item \textbf{Модели хрупкости} (гамма-распределённый индивидуальный множитель риска). Учитывают отбор
        по смертности --- хрупкие умирают первыми, поэтому выжившие представляют собой отобранную
        подгруппу --- и являются общепризнанной поправкой. Не используются, потому что распределение
        хрупкости неидентифицируемо по агрегированным числам смертей, и оно молча изменило бы смысл каждого
        возрастного показателя.
  \item \textbf{Явные конкурирующие риски с копулой.} Статистически наиболее полный подход, допускающий
        зависимость двух причин смерти; требует информации о степени зависимости, которую агрегированные
        поперечные данные дать не могут.
\end{itemize}}

\section{Возникновение рака: риск диагностики}

\subsection{Функциональная форма риска}

\mech{Годовая вероятность того, что рак будет диагностирован клинически в возрасте $a$ (при условии, что
диагноза не было раньше), --- это \emph{кусочно-постоянный риск на логарифмической шкале}:
\[
\log h(a) \;=\; \beta_0 + \sum_{k=1}^{K}\beta_k\,\mathbf{1}\{a \ge t_k\},
\]
где возрастные точки разрыва $t_k$ и коэффициенты $\beta_k$ задаются конфигурацией (в комплекте по умолчанию
$t = (-1,0,40,50,60,70)$ и $\beta = (-11.9,0.01,0.28,-0.11,-0.1,-0.1)$). Значение $h(a)$ ограничено сверху
единицей. \emph{Возраст диагностики} агента --- это первый возраст $a$, при котором испытание Бернулли с
вероятностью $h(a)$ даёт успех, при просмотре возрастов начиная с нуля; если ни один возраст не дал успеха,
агент никогда не диагностируется и получает служебный возраст диагностики ($110$ --- «нереальная
продолжительность жизни»), означающий «клинического рака нет». Просмотр векторизован: для всех агентов
столбец равномерных величин сравнивается с кривой риска и берётся первое попадание.}

\why{Функциональная форма прямо следует из того, что можно откалибровать и что нужно модели дальше.

\emph{(i) Цель калибровки --- возрастной годовой показатель заболеваемости}, а кусочно-постоянный
дискретный риск --- это ровно та величина, которая ему соответствует: «годовая вероятность диагноза при
условии, что диагноза ещё не было». Выбор риска (а не распределения возраста диагностики) делает шаг
подгонки линейным по коэффициентам, и он решается за миллисекунды, что и обеспечивает интерактивный
веб-интерфейс.

\emph{(ii) Точки разрыва поставлены там, где наблюдаемая кривая меняет наклон.} Заболеваемость круто
растёт с середины жизни, а узлы 40, 50, 60 и 70 --- это возрасты смены наклона в целевых данных, поэтому
несколько параметров описывают форму, не навязывая функциональное семейство.

\emph{(iii) Логарифмическая шкала делает две вещи сразу.} Она сохраняет $h$ строго положительной без
ограничений и позволяет константе выражать базовую вероятность в нескольких порядках величины:
подогнанное $\beta_0 \approx -11.9$ соответствует риску порядка $10^{-5}$ в год в молодых возрастах, что
является правильным порядком для редкого рака и что неудобно представлять на линейной шкале.

\emph{(iv) Ограничение сверху единицей и выбор первого успеха} превращают величину в настоящую вероятность
и в настоящее время первого события. Ограничение необходимо, потому что механизму Бернулли нужна
вероятность; выбор первого успеха кодирует тот факт, что «возраст диагностики» --- это время первого
события, поэтому ни один агент не может быть диагностирован дважды и не может накопить два независимых
диагноза.}

\alt{\begin{itemize}[nosep]
  \item \textbf{Гладкий параметрический риск} (Гомпертц, Вейбулл, логнормальный, гамма). Один-два
        параметра, гладкая производная, просто для изложения; риск --- систематическое расхождение именно
        там, где это важно: на перегибе и на позднем плато заболеваемости. Естественный выбор, если цель ---
        один узкий возрастной диапазон, а не вся жизнь.
  \item \textbf{Сплайны} (натуральные кубические, B-сплайны, штрафные P-сплайны). Гибче ступенчатых функций
        и глаже кусочно-постоянной подгонки; цена --- дополнительные параметры и выбираемый вес штрафа. Это
        наиболее естественное развитие текущей конструкции.
  \item \textbf{Использование наблюдаемых возрастных показателей напрямую как риска} (без подгонки вовсе).
        Полностью снимает вопрос идентифицируемости; цена --- выборочный шум, отсутствие сглаживания,
        невозможность интерполяции и экстраполяции и тонкое расхождение между годовым показателем и годовым
        дискретным риском без поправки на человеко-годы.
  \item \textbf{Механистическая модель канцерогенеза} (двухстадийная клональная экспансия
        Мулгкавкара --- Вензона --- Кнудсона, многостадийная модель Армитеджа --- Долла, многостадийная с
        геномной нестабильностью). Объясняет, \emph{почему} заболеваемость растёт примерно как степень
        возраста, и связывает инициацию и промоцию с экспозициями, что сделало бы возможным моделирование
        факторов риска. Не используется, потому что такие модели неидентифицируемы по одной возрастной
        заболеваемости --- множество комбинаций параметров воспроизводят ту же кривую, --- и потому что
        дополнительная структура не улучшает ответ на вопрос о скрининге.
  \item \textbf{Возраст-период-когортное моделирование.} Правильный инструмент для \emph{трендов}
        заболеваемости; ему нужны несколько периодов и когорт, тогда как эта модель сознательно
        моделирует один поперечный срез.
  \item \textbf{Раздельные риски для прогрессирующих и непрогрессирующих образований} ---
        последовательность «аденома--карцинома» при колоректальном раке с собственной клинически
        обнаружимой доклинической фазой. Это структурно правильная трактовка предрака, и это самое крупное
        структурное упущение текущей модели (раздел~\ref{sec:assumptions}); оно требует данных уровня
        образований.
\end{itemize}}

\subsection{Что считается «наличием рака», и самые старые агенты}

\mech{Агент помечается \id{has\_cancer} тогда и только тогда, когда разыгранный возраст диагностики ниже
служебного значения ($110$) \emph{и} ниже естественного возраста смерти агента. Человеку, который умер бы
раньше года клинической диагностики, история опухоли не создаётся вовсе. Кроме того, агенты, которым на
$t=0$ уже больше 98 лет, сразу получают служебное значение, поэтому в ходе прогона заболеть раком они не
могут.}

\why{Оба правила следуют из одного решения: единица рака в модели --- это \emph{рак, который был бы
клинически диагностирован в течение оставшейся жизни данного человека}. Именно это событие порождает
наблюдаемые заболеваемость и смертность и именно то событие, которое способен сместить скрининг. Агент,
умирающий до года диагностики, не вносит вклада ни в один отчётный показатель, поэтому построение для него
опухолевой истории добавило бы стоимость и бессмысленную траекторию. У возрастного ограничения тот же
мотив в более грубой форме: просмотр возрастов диагностики ограничен максимальным представимым возрастом, а
обратный расчёт возраста начала доклинической фазы для агента старше 98 лет поместил бы всю доклиническую
траекторию за пределы моделируемого окна, поэтому такие агенты просто исключаются.}

\why{\emph{Чего это стоит, сказано прямо.} Естественная история \emph{управляется диагностикой}: у всех,
у кого в модели есть рак, он был бы диагностирован клинически. Поэтому модель не может представить рак,
который существует, но никогда не становится клинически явным, --- а это и есть механизм гипердиагностики.
Без скрининга это безвредно, поскольку такой рак всё равно невидим ни в одном отчётном показателе. При
скрининге это уже не безвредно: каждый обнаруженный при скрининге рак --- «настоящий» рак, который был бы
диагностирован и мог бы убить, поэтому при наличии реальной доли непрогрессирующих образований модель
завышает пользу скрининга. Возрастное ограничение --- тоже грубое упрощение без эпидемиологического
обоснования: это диапазонный предохранитель, оформленный как правило риска, --- и его следует пересмотреть,
если изменятся горизонт, начальная возрастная структура или служебный возраст.}

\alt{\begin{itemize}[nosep]
  \item \textbf{Естественная история, управляемая началом заболевания.} Выдать каждому агенту доклиническое
        начало (по риску инициации) и позволить прогрессии, симптомам и скринингу определять, будет ли рак
        вообще и когда диагностирован. Тогда человек может умереть с недиагностированным раком, и как
        гипердиагностика, так и «обнаружено при скрининге, но никогда бы не проявилось» становятся
        представимыми. Это структурно правильная конструкция и направление серьёзного улучшения; она требует
        модели прогрессии и данных о доклинической обнаруживаемости, которых в комплекте нет.
  \item \textbf{Доля непрогрессирующих (или медленно прогрессирующих) образований.} Стандартный приём в
        моделях скрининга: доля $\mathrm{nc}$ образований никогда не достигает клинической значимости. Один
        дополнительный параметр, и это принятый способ сделать гипердиагностику измеримой --- но его нельзя
        оценить только по заболеваемости и смертности, поэтому его пришлось бы задать из литературы и
        сделать объектом анализа чувствительности.
  \item \textbf{Никаких исключений: строить опухолевую историю для каждого агента независимо от его судьбы.}
        Концептуально единообразно и устраняет особый случай, но умножает работу ради информации, которую
        не использует ни один выходной показатель, и порождает «рак у уже умерших».
  \item \textbf{Убрать ограничение 98 лет, расширив служебное значение.} Чисто техническое исправление:
        ограничение следует либо снять (пусть решает риск), либо заменить явно задокументированным
        возрастом усечения.
\end{itemize}}

\section{Доклиническая фаза и проблема масштаба времени}

\subsection{Экспоненциальное время до диагностики}

\mech{Для каждого агента с раком модель разыгрывает \emph{время до диагностики} (длительность доклинической
фазы, то есть время от первой злокачественной клетки до клинического диагноза) из экспоненциального
распределения со средним \id{lead\_time\_mean} (3,5 года в конфигурации по умолчанию), то есть
$t_{\text{ttd}} \sim \mathrm{Exp}(\lambda)$, $\lambda = 1/\id{lead\_time\_mean}$. Затем возраст начала
заболевания вычисляется обратным счётом от \emph{наблюдаемого} возраста диагностики:
$\text{начало} = \max(\text{диагноз} - \lceil t_{\text{ttd}}\rceil, 0)$. Непрерывное значение
$t_{\text{ttd}}$ используется для определения стадии (раздел~\ref{sec:stage}), а целое --- для арифметики
возрастов.}

\why{Вся доклиническая история привязывается к дате диагностики, а не моделируется вперёд от даты начала, и
именно это решение делает возможной калибровку по агрегированным данным.

\emph{(i) Наблюдаемая величина --- это возраст диагностики.} Привязка доклинической траектории к нему
означает, что модель воспроизводит наблюдаемую заболеваемость по построению, а время до диагностики
используется лишь для того, чтобы определить, как выглядела опухоль \emph{до} уже правильной даты. Ничего
не должно быть верным в абсолютной скорости начала заболевания, поэтому целый класс расхождений устраняется.

\emph{(ii) Экспоненциальное распределение --- выбор с максимальной энтропией при заданном среднем.}
Предположение постоянной скорости прогрессии (без памяти о том, как долго опухоль уже существует) ---
наименее обязывающее из доступных для времени пребывания, и именно его использует классическая марковская
литература о скрининге. Такая совместимость не косметична: соотношения между временем до диагностики,
временем пребывания и пользой скрининга, выведенные в этой литературе, остаются применимыми для проверки
поведения модели.

\emph{(iii) Один параметр, потому что позволить можно только один.} Абсолютный масштаб доклинической фазы
неидентифицируем (раздел~\ref{sec:nonident}), поэтому второй параметр времени пребывания добавил бы вторую
неопределяемую величину. Среднее --- наименьшая достаточная сводка: это масштаб всего, что ниже по цепочке
(времена пересечения, возраст начала, длительность стадии), и единственная ручка, которую поворачивает
анализ чувствительности.

\emph{(iv) Отсутствие памяти сохраняет калибровку стадий в замкнутой форме.} При экспоненциальном времени
до диагностики $P(\text{стадия} \ge s) = \exp(-t_s/\mu)$, и именно это позволяет \emph{обратить} наблюдаемое
распределение по стадиям во времена пересечения одним логарифмом на стадию (раздел~\ref{sec:stage}), а не
решать численную задачу калибровки.

\emph{(v) Непрерывная величина для сравнения, целая для арифметики.} Округление времени до диагностики до
сравнения с временами пересечения сместило бы распределение стадий (округление вниз систематически
порождало бы больше ранних стадий), поэтому сравнение использует непрерывную величину, а арифметика
возрастов --- её потолок.}

\alt{\begin{itemize}[nosep]
  \item \textbf{Гамма- или вейбулловское время пребывания.} Стандартные обобщения, широко применяемые в
        моделях обнаружения при скрининге; экспоненциальное --- их частный случай при параметре формы 1. Они
        дают более реалистичную форму (несколько очень долгих пребываний), но параметр формы приходится
        задавать, а не оценивать, то есть одно предположение обменивается на два.
  \item \textbf{Детерминированное (фиксированное) время до диагностики.} Простейшее предположение; тогда
        распределение по стадиям становится функцией одного порога «по острию ножа», а не гладкой
        вероятностью, и доклиническая длительность становится нереалистично дискретной.
  \item \textbf{Длительности на каждую стадию} (время пребывания для каждой стадии). Гораздо ближе к
        биологии опухоли и делает длительность стадии прямо интерпретируемой; доступные данные не могут
        обосновать четыре отдельные длительности, и это заставляет времена пересечения лежать на целочисленной
        сетке.
  \item \textbf{Обнаруживаемость по размеру, выводимая из самой кривой роста.} Если опухоль «обнаружима»
        после пересечения порога размера, время пребывания --- это то, что говорит модель Гомпертца, что
        красиво самосогласованно, поскольку эта кривая в модели уже есть. Не хватает порогового размера
        обнаруживаемости и зависящей от размера чувствительности теста, которых в комплекте данных нет.
  \item \textbf{Оценка времени до диагностики по данным скрининга} (обнаруженные при скрининге против
        интервальных раков, или данные рандомизированного исследования скрининга). Это единственный путь к
        подлинной идентификации и самая ценная добавка данных, какую модель могла бы получить.
  \item \textbf{Привязка среднего к литературе} вместо свободного числа 3,5, с опубликованным диапазоном в
        качестве интервала анализа чувствительности. Операционно идентично текущей конструкции, но
        превращает предположение в цитируемое предположение --- самое дешёвое из доступных улучшений.
\end{itemize}}

\subsection{Неидентифицируемость масштаба времени и что вместо него фиксируется}
\label{sec:nonident}

\mech{Модель \emph{не} оценивает масштаб доклинического времени. Она фиксирует среднее время до диагностики
(\id{lead\_time\_mean}) и калибрует всё остальное, а времена пересечения стадий отдельно \emph{выводит} из
наблюдаемого распределения по стадиям (раздел~\ref{sec:stage}). Зафиксированный масштаб затем предъявляется
как единственная ось анализа чувствительности (раздел~\ref{sec:sens}).}

\why{Причина --- подлинное математическое свойство системы, а не недостаток усилий. Обозначим
$c_s = C - \ln(K - \ln s)$ и запишем время пересечения стадии $s$ как $t_s = c_s / B_{\text{ind}}$
(раздел~\ref{sec:gompertz}). При экспоненциальном времени до диагностики со средним $\mu$
\[
P(\text{диагноз на стадии} \ge s) \;=\; P(t_s \le t_{\text{ttd}}) \;=\; \exp\!\left(-\frac{c_s}{B_{\text{ind}}\,\mu}\right),
\]
то есть наблюдаемое распределение по стадиям ограничивает только \emph{произведение}
$B_{\text{ind}}\mu$. Удвоение скорости роста при вдвое меньшем среднем времени пребывания оставляет все
наблюдаемые величины ровно такими же. Формально вектор параметров $(B_{\text{ind}}, \mu)$ не идентифицируем
вдоль гиперболы $B_{\text{ind}}\mu = \text{const}$ --- это плоский гребень функции правдоподобия, и
максимизатор правдоподобия охотно вернул бы произвольную точку на нём.

Есть и вторая, более фундаментальная причина. Начало заболевания --- появление первой злокачественной
клетки --- не является наблюдаемым событием. Поэтому время до диагностики --- это не величина, которая
плохо измеряется; это величина, в принципе ненаблюдаемая в клиническом наборе данных, поскольку диагностика
прерывает именно ту историю, которую хотелось бы наблюдать. «Оценённое время до диагностики» в литературе
всегда означает «оценённое в рамках явной модели».

Что данные \emph{действительно} идентифицируют: риск диагностики (по возрастной заболеваемости),
произведение $B_{\text{ind}}\mu$ (по распределению стадий) и риск терминальной фазы (по смертности от рака).
Поэтому конструктивное решение --- потратить оставшуюся степень свободы как \emph{явное предположение с
интерпретируемым именем}, а не как подогнанное значение, притворяющееся оценкой. Фиксация именно $\mu$, а не
$B$, --- асимметричный, но правильный выбор: для колоректального рака у $\mu$ есть оценки в литературе и
простой клинический смысл (время от начала заболевания до диагноза), тогда как масштаб $B$ не с чем
соотнести вне модели.

Практическое следствие --- чистое разделение труда: $\mu$ есть предположение, времена пересечения
\emph{выводятся} из данных, а параметры скорости роста \emph{подгоняются} под эти выведенные времена.
Поэтому модель воспроизводит наблюдаемое распределение стадий по построению, а единственное предположение
видно и в файле параметров, и в документации, и на графиках чувствительности.}

\alt{\begin{itemize}[nosep]
  \item \textbf{Сделать диапазон времени до диагностики заголовочным результатом, а не приложением.} В
        математике ничего не меняется; меняется то, что читатель видит как ответ семейство кривых, а не одну
        кривую с оговоркой. Для вопроса о скрининге это, вероятно, более честная подача, и это изменение
        отчётности, а не модели.
  \item \textbf{Байесовский подход: априорное распределение на $\mu$, остальное подгоняется, полное
        апостериорное распределение.} Статистически самый чистый способ сказать «$\mu$ есть суждение» ---
        суждение входит как априорное распределение, а его последствия приходят как доверительные интервалы
        для всех выходных показателей, а не как набор повторных прогонов. Требует MCMC (или ABC) ценой часов
        вместо миллисекунд и делает заголовочные числа зависящими от априорного распределения, которое само
        не выведено из данных.
  \item \textbf{Профильное правдоподобие по $\mu$.} Вычислить качество подгонки как функцию $\mu$ и показать
        поверхность. Дёшево (модель быстра), прозрачно и делает гребень видимым, а не только описанным. Это
        минимальный возможный шаг к корректному учёту неопределённости.
  \item \textbf{Зафиксировать $B$ вместо $\mu$.} Симметрично и столь же произвольно, и однозначно хуже: у
        скорости роста вообще нет внешней привязки, поэтому выбор был бы неинтерпретируемым и
        непроверяемым.
  \item \textbf{Время пребывания с двумя параметрами.} Хуже, а не лучше: один неопределяемый масштаб
        заменяется двумя неопределяемыми параметрами, и никакое количество экспертного мнения их не
        идентифицирует.
  \item \textbf{Внешняя идентификация по данным скрининга.} Исследование скрининга или регистр обнаруженных
        при скрининге и интервальных раков идентифицируют время пребывания напрямую; это единственный способ
        превратить $\mu$ в оценку. Требуются индивидуальные данные скрининга, которых у модели нет.
\end{itemize}}

\section{Рост опухоли: редуцированная модель Гомпертца}
\label{sec:gompertz}

\subsection{Кривая роста}

\mech{Рост опухоли описывается \emph{редуцированной функцией Гомпертца}
\[
V(t) = \exp\!\bigl(K - \exp(C - B\,t)\bigr),
\]
где $V$ --- абстрактный суррогат массы опухоли («размер»), $t$ --- годы с момента появления первой
злокачественной клетки, $K$ и $C$ --- параметры формы, $B$ --- скорость роста. Кривая начинается со
значения $V(0) = \exp(K - \exp C)$ (небольшая опухоль), насыщается на $V(\infty) = \exp K$ (ёмкость среды) и
имеет точку перегиба при $t = C/B$, где рост переходит от ускорения к замедлению. Стадии --- это времена
пересечения размеров $2$, $3$ и $4$:
\[
t_s = \frac{C - \ln(K - \ln s)}{B}, \qquad s \in \{2,3,4\},
\]
конечное ровно тогда, когда $\exp K > s$, то есть когда ёмкость среды выше порога стадии. В эталонной
подгонке ($K = 1.451$, $C = 0.198$, $B = 0.550$, так что $\exp K = 4.27$) времена пересечения равны
$0.87$, $2.26$ и $5.34$ года.}

\why{Гомпертц выбран по четырём причинам: три содержательные и одна структурная.

\emph{(i) Это классический закон роста опухоли, и он соответствует качественной форме данных.} Сначала
опухоли растут примерно экспоненциально, затем замедляются, исчерпывая кровоснабжение и доступное
пространство; кривая Гомпертца обладает ровно этой особенностью --- экспоненциальна на раннем этапе и
асимптотически плоская на позднем --- и полвека является описательной моделью роста по умолчанию в
онкологии.

\emph{(ii) Он экономичен, но не жёсток.} Два параметра формы ($K$, $C$) и одна скорость ($B$) воспроизводят
широкий диапазон кривых роста, и у каждого параметра есть ясный смысл: $\exp K$ --- ёмкость среды в
единицах размера, $C$ задаёт начальный размер через $V(0) = \exp(K - \exp C)$, а $B$ --- удельная скорость
роста.

\emph{(iii) Его времена пересечения имеют замкнутую форму.} Обратная функция
$t_s = (C - \ln(K - \ln s))/B$ делает калибровку стадий задачей, \emph{линейной относительно логарифма}, и
позволяет сохранить возрасты стадий каждого агента на этапе генерации. Закон роста, чьи времена стадий
требовали бы численного обращения, был бы и медленнее, и гораздо труднее для рассуждений --- а замкнутая
форма ещё и позволяет доказать утверждение о неидентифицируемости из раздела~\ref{sec:nonident}, а не
установить его эмпирически.

\emph{(iv) Пересечения абсолютных размеров $\{2,3,4\}$ --- это нормировка, а не ограничение.} В модели три
времени пересечения как цели калибровки и три свободных параметра ($K$, $C$, $B$), поэтому подгонка
идентифицирована точно: соглашение о порогах поглощается подобранными формами. На что расстояние между
порогами \emph{действительно} влияет --- так это на индивидуальный разброс стадий (через
$B_{\text{std}}$, раздел~\ref{sec:hetero}), поскольку разные опухоли пересекают по-разному расположенные
пороги в разные моменты. Это ограничение, которое нужно держать в виду, а не скрытое искажение
популяционных средних.}

\alt{\begin{itemize}[nosep]
  \item \textbf{Экспоненциальный рост.} Простейший возможный закон; неограниченный, поэтому без ёмкости
        среды, и он заставляет времена пересечения быть равноотстоящими по логарифму размера --- сильное
        ограничение, которое данные опровергли бы.
  \item \textbf{Логистический рост.} Знаком и симметричен; замедление симметрично относительно точки
        перегиба, что для опухолей, как известно, неверно (они долго растут, прежде чем замедлиться), и по
        той же причине симметричны времена пересечения.
  \item \textbf{Фон Берталанфи / Берталанфи --- Пюттер.} Механистически обоснован (рост как баланс
        анаболических и катаболических слагаемых) и принадлежит тому же семейству, что и Гомпертц;
        сопоставимое поведение, менее привычен в литературе о скрининге.
  \item \textbf{Вейбулловский или степенной рост.} Гибок; нет ёмкости среды и нет устоявшейся клинической
        интерпретации параметров.
  \item \textbf{Механистическая модель, ограниченная питанием или пространством, с явной васкуляризацией.}
        Наиболее биологически обоснованная опция; много параметров и нужны объёмные данные о росте, которых
        агрегированная онкологическая статистика не содержит.
  \item \textbf{Стохастический ветвящийся рост.} Моделирует клональную динамику напрямую и --- уникально ---
        способен выражать \emph{вымирание} ранних опухолей, а это настоящий механизм непрогрессии.
        Вычислительно гораздо тяжелее и неидентифицируем по доступным данным.
  \item \textbf{Кусочно-линейная прогрессия (без модели роста).} Постоянная скорость прогрессии для класса
        опухолей, то есть классическая многосостоянийная марковская формулировка. Проще и подгоняется прямо
        под данные о стадиях; она разрывает связь между ростом и сроками стадий, которая и делает текущую
        модель интерпретируемой, и оставляет понятие «размер» без всякого содержания.
  \item \textbf{Прямая порядковая регрессия стадии по ковариатам} --- конструкция, которую это заменило. Она
        воспроизводит наблюдаемое распределение стадий, но стадия становится произвольным разыгрыванием,
        оторванным от процесса роста, индивидуальная неоднородность теряет физический смысл, а ошибка
        остаётся незаметной, потому что результаты выглядят правдоподобно.
\end{itemize}}

\subsection{Индивидуальная неоднородность скорости роста}
\label{sec:hetero}

\mech{Каждый агент с раком получает собственную скорость роста
\[
B_{\text{ind}} = \exp\!\bigl(\ln B_{\text{pop}} + \varepsilon\bigr), \qquad
\varepsilon \sim \mathcal{N}(0, B_{\text{std}}^2),
\]
так что $B_{\text{ind}}$ логнормально распределена с медианой $B_{\text{pop}}$ и мультипликативным разбросом
$\exp(B_{\text{std}})$ (нижняя граница $10^{-6}$ исключает вырожденный ноль). Кривая популяционного среднего,
используемая для подгонки, берётся при $B_{\text{pop}}$ без случайного эффекта. $B_{\text{std}}$
\emph{не} подгоняется: это заданное вручную значение (по умолчанию 0,1, то есть мультипликативный разброс
примерно $\pm10\%$).}

\why{Два аргумента: один о реалистичности и один о конкретном феномене скрининга.

\emph{(i) Скорости роста действительно различаются, и этот клинический факт --- не мешающий параметр.}
Время удвоения опухоли при одном и том же раке различается между пациентами в разы, а не на проценты. Без
неоднородности стадия является детерминированной функцией времени до диагностики: каждый пациент с тем же
истекшим доклиническим временем оказался бы на той же стадии, что навязывает жёсткую и искусственную
корреляцию между «как долго растёт опухоль» и «насколько она запущена». Неоднородность их разводит, поэтому
два пациента с одинаковым временем до диагностики вполне законно могут быть диагностированы на разных
стадиях.

\emph{(ii) Смещение по длительности требует именно её.} Скрининг предпочтительно выявляет медленно растущие
опухоли, потому что опухоль, проводящая больше времени в обнаруживаемой доклинической фазе, имеет больше
шансов быть пойманной. Это смещение по длительности (length bias) --- вместе со связанным с ним смещением
по времени упреждения --- одна из классических причин, по которым исследования скрининга переоценивают
пользу, и оно может появиться в модели только тогда, когда скорости роста различаются между людьми. При
$B_{\text{std}} = 0$ этот феномен исчезает по построению, и сдвиг стадий в модели становится слишком
благоприятным и слишком «чистым».

Логнормальное распределение, а не нормальное, --- естественная параметризация для скорости: положительность
автоматична, распределение правоскошенное (несколько быстрых опухолей, «пол» медленных), как это обычно и
бывает у мультипликативных величин, а оба параметра имеют прямой смысл --- $B_{\text{pop}}$ есть медиана, а
$B_{\text{std}}$ --- разброс по логарифмической шкале. Одна тонкость, которую стоит помнить: среднее
арифметическое равно $B_{\text{pop}}\exp(B_{\text{std}}^2/2)$, то есть чуть выше медианы, поэтому сравнение
на основе среднего с единственной подобранной скоростью несёт небольшое, но известное смещение.}

\why{\emph{Честно названная слабость.} $B_{\text{std}}$ определяет, сколько индивидуальной вариации,
разброса стадий и смещения по длительности выдаёт модель, --- и он задан вручную, а не оценён. Ничто в
текущей калибровке его не ограничивает, поскольку шаг подгонки нацелен только на времена пересечения
популяционного среднего ($K$, $C$, $B_{\text{pop}}$). Это, следовательно, второе непроверенное
предположение наряду со средним временем до диагностики, но, в отличие от последнего, оно сейчас
\emph{не} варьируется в анализе чувствительности. Превратить его в ось анализа чувствительности дёшево и
существенно усилило бы утверждение о неопределённости.}

\alt{\begin{itemize}[nosep]
  \item \textbf{Без неоднородности ($B_{\text{std}} = 0$).} Простейшая опция, при которой стадия становится
        детерминированной функцией времени до диагностики. Она полностью убирает смещение по длительности и
        разброс стадий и делает корреляцию между стадией, временем до диагностики и выживаемостью
        искусственно жёсткой, --- поэтому её следует использовать только как осознанное сравнение в анализе
        чувствительности, но не как базовый вариант.
  \item \textbf{Гамма-распределённая скорость роста.} Тоже положительна, с независимым управлением формой и
        масштабом; менее естественна для мультипликативной величины и не лучше соответствует данным.
  \item \textbf{Несколько дискретных классов роста} (медленный / средний / быстрый). Очень просто для
        изложения и для сценарного анализа; грубо, вносит искусственную дискретность, а границы классов
        произвольны.
  \item \textbf{Нормальная (неусечённая) скорость роста.} Отвергнуто: допускает отрицательные скорости,
        которые не просто маловероятны, а бессмысленны.
  \item \textbf{Полная модель случайных эффектов для $\ln B$ с ковариатами} (возраст, пол, стадия в начале).
        Правильный способ сделать неоднородность роста явной, а не побочной; нужна информация о росте с
        разбивкой по ковариатам, которой в файле стадий нет.
  \item \textbf{Эмпирические времена удвоения объёма} из серий визуализации или из литературы. Единственный
        путь к защитимому значению $B_{\text{std}}$; оно превратило бы заданное вручную число в цитируемое.
\end{itemize}}

\subsection{От роста к стадии: стадия как выводимая величина}
\label{sec:stage}

\mech{Для каждого агента с раком модель вычисляет его собственные времена пересечения $t_2, t_3, t_4$ (по
$B_{\text{ind}}$, раздел~\ref{sec:gompertz}) и сравнивает их с непрерывным временем до диагностики: стадия
при клинической диагностике равна $4$, если $t_4 \le t_{\text{ttd}}$, иначе $3$, если
$t_3 \le t_{\text{ttd}}$, иначе $2$, если $t_2 \le t_{\text{ttd}}$, иначе $1$. Иными словами, стадия ---
это \emph{число порогов размера, которые опухоль пересекла к моменту клинической диагностики}. Полный массив
возрастов стадий, сохраняемый на агента, есть
$[\text{начало},\ \text{начало}+t_2,\ \text{начало}+t_3,\ \text{начало}+t_4,\ \text{начало}+t_4+\delta]$,
где $\delta$ --- экспоненциальная величина для терминальной фазы после последнего пересечения.

Времена пересечения, используемые как \emph{цели} калибровки, получаются обращением наблюдаемого
распределения по стадиям. Поскольку время до диагностики экспоненциально,
$P(\text{стадия} \ge s) = \exp(-t_s/\mu)$, поэтому
\[
t_2 = -\mu \ln(1 - p_1), \qquad t_3 = -\mu \ln(p_3 + p_4), \qquad t_4 = -\mu \ln p_4 ,
\]
где $p_s$ --- наблюдаемая доля стадии $s$, а $\mu$ --- среднее время до диагностики. Для файла стадий
колоректального рака по США наблюдаемое распределение равно $0.220/0.256/0.308/0.218$, что даёт цели
$0.87/2.26/5.34$ года; полученная подгонка воспроизводит распределение как
$0.225/0.250/0.284/0.241$ ($N = 20\,000$, зерно 2022).}

\why{Стадия выводится, а не разыгрывается, и разница значит больше, чем кажется.

\emph{(i) Это делает распределение стадий следствием, а не входом.} Когда стадия разыгрывается из
распределения, воспроизведение наблюдаемого распределения --- тавтология: модели сообщили ответ. Здесь
распределение порождается взаимодействием двух физических величин (как долго существует опухоль и как быстро
она растёт), поэтому согласие с наблюдаемым распределением становится настоящей, опровержимой проверкой
калибровки. Остаточное расхождение ($0.225/0.250/0.284/0.241$ против $0.220/0.256/0.308/0.218$) мало и
объясняется выборочным шумом, $B_{\text{std}}$ и взаимодействием стадии с зависящим от возраста риском
диагностики --- но не систематическим несоответствием.

\emph{(ii) Именно это придаёт смысл индивидуальной скорости роста.} Если бы стадия происходила из
регрессии, $B_{\text{ind}}$ было бы числом без последствий; поскольку времена пересечения вычисляются из
него, скорость роста управляет стадией, возрастами стадий и, следовательно, сроками терминальной фазы.
Неоднородность роста имеет наблюдаемые последствия --- ровно то, что и должно быть у структурного параметра.

\emph{(iii) Обращение имеет замкнутую форму, поэтому цель калибровки выводится из данных и проверяема.} Один
логарифм на стадию превращает файл стадий в три числа, которые можно напечатать, проверить вручную и
проследить до исходной таблицы. Нет внутреннего цикла оптимизации, в котором могла бы спрятаться ошибка.

\emph{(iv) Это согласуется с учётом неидентифицируемости.} Обращение потребляет $\mu$ (принятый масштаб) и
возвращает времена пересечения; подгонка Гомпертца затем потребляет времена пересечения и возвращает $K$,
$C$, $B_{\text{pop}}$. Три цели, три параметра --- точно идентифицированная цепочка, в которой единственное
предположение видно в самом её начале.}

\alt{\begin{itemize}[nosep]
  \item \textbf{Разыгрывать стадию прямо из наблюдаемого распределения.} Точно воспроизводит распределение
        ценой того, что стадия становится экзогенной: неоднородность роста, время до диагностики и скрининг
        уже не могут сдвинуть стадию иначе как явной произвольной подгонкой, поэтому модель роста опухоли
        становится декоративной.
  \item \textbf{Многочленная или порядковая регрессия стадии по возрасту, полу и другим ковариатам.}
        Конструкция, которую это заменило. Она хорошо подгоняет распределение и проста в реализации, но
        отрывает стадию от процесса роста, оставляет индивидуальную скорость роста бессмысленной, а
        интерпретация коэффициентов («влияние возраста на логарифм шансов более поздней стадии») не имеет за
        собой механизма.
  \item \textbf{Латентный процесс прогрессии по стадиям со скоростями переходов} (стандартная
        многосостоянийная формулировка). Чисто и делает длительности явными; нужны наблюдаемые данные о
        переходах между стадиями (например, из последовательных раундов скрининга), которых в комплекте нет.
  \item \textbf{Клиническое стадирование TNM} (размер опухоли $T$, поражение узлов $N$, отдалённые
        метастазы $M$). Клинически правильное и гораздо более информативное описание; позволило бы измерять
        реальный сдвиг стадий при скрининге относительно настоящих порогов, но требует информации о размере
        и узлах, которой файл со стадиями не содержит.
  \item \textbf{Дискретные классы роста с фиксированными длительностями стадий} (медленные, промежуточные,
        быстрые опухоли, каждая с фиксированным временем пребывания на стадии). Просто и прозрачно, и именно
        так параметризованы некоторые опубликованные модели скрининга; разброс стадий становится дискретным
        выбором, а не непрерывным.
\end{itemize}}

\section{Агрессивность: метка опухоли, изменяющая шансы излечения}
\label{sec:agg}

\subsection{Что это за метка и на что она влияет}

\mech{Каждый агент с раком помечается как \emph{агрессивный} или нет испытанием Бернулли, вероятность
которого берётся из файла стадий: для опухоли, диагностированной на стадии $s$, вероятность равна
наблюдаемой доле агрессивных опухолей на этой стадии, вычисленной по колонке \id{Aggressiveness} (данные по
колоректальному раку США: $0.096 / 0.198 / 0.437 / 0.639$ для стадий I--IV, то есть агрессивность и стадия
растут вместе). Затем метка влияет \emph{ровно на одно} --- на вероятность излечения. Базовая вероятность
излечения $p$ переводится в шансы, умножается на отношение шансов для данной стадии
(\id{aggressiveness\_cure\_odds\_ratio}, по умолчанию $0.667$ для всех стадий) и переводится обратно:
\[
\text{шансы}' = \frac{p}{1-p}\times \mathrm{OR}_s, \qquad
p' = \frac{\text{шансы}'}{1 + \text{шансы}'}.
\]
Значение $\mathrm{OR}_s = 1$ строго ничего не меняет (преобразование пропускается), поэтому предположение
можно отключить, не затрагивая никакого другого поведения.}

\why{Четыре причины, и они связаны между собой.

\emph{(i) Это единственное измерение фенотипа, по которому у модели есть данные.} В файле стадий есть
признак агрессивности, поэтому доля по стадиям \emph{измерена}, а не принята на веру. Игнорировать его
означало бы выбросить реальную информацию; придумать другую ось фенотипа опухоли означало бы не иметь за
ней никаких данных вообще.

\emph{(ii) Он сознательно не является модификатором роста, чтобы избежать двойного учёта.} Агрессивные
опухоли, по клинической интуиции, растут быстро, поэтому соблазнительная конструкция --- позволить метке
менять $B$. Это смешало бы её с $B_{\text{std}}$, который уже несёт неоднородность роста, и калиброванное
распределение стадий должно было бы поглотить два перекрывающихся источника вариации. Отнесение метки к
вероятности излечения помещает её туда, где она наблюдаема: в выживаемость, а не в стадирование.

\emph{(iii) Форма отношения шансов --- правильная алгебра для вероятности излечения.} Излечение есть бинарный
исход, поэтому отношение шансов является стандартной мерой эффекта; оно мультипликативно сочетается с
базовыми шансами, не может выйти за пределы $[0,1]$, симметрично по формулировке ($\mathrm{OR}$ и
$1/\mathrm{OR}$ выражают «труднее» и «легче» одним и тем же параметром) и обращается в тождество при
$\mathrm{OR} = 1$.

\emph{(iv) Направление --- клинический априорный довод, а форма оставляет место для данных по стадиям.}
То, что агрессивную опухоль труднее излечить, --- разумный априорный довод, и $\mathrm{OR}<1$ его выражает.
Параметр является вектором по стадиям, хотя по умолчанию все четыре равны, поэтому пользователь, у которого
есть данные об ответе на лечение по стадиям, может их ввести, не меняя конструкцию.}

\why{\emph{Две ограничения, которые нужно назвать прямо.}
Во-первых, само отношение шансов \emph{не оценено по данным} --- оно задано вручную. В файле стадий нет
колонки об исходе лечения, поэтому в комплекте входных данных нет ничего, что могло бы его
идентифицировать: то, что измеримо (доля агрессивных по стадиям), измеряется, а то, что неизмеримо (влияние
метки на излечение), является явным предположением с клиническим априорным доводом.
Во-вторых, поскольку агрессивность сильно связана со стадией ($0.096$ против $0.639$), объединённое
сравнение выживаемости агрессивных и неагрессивных смешивает отношение шансов с составом по стадиям ---
агрегированные числа покажут большой разрыв в выживаемости, который в основном является эффектом стадии.
Это ловушка парадокса Симпсона, и модель её избегает, представляя кривые выживаемости и доли излечения
\emph{с разбивкой по стадии и агрессивности}, а не только объединённо (раздел~\ref{sec:survival}).}

\alt{\begin{itemize}[nosep]
  \item \textbf{Агрессивность, влияющая на рост ($B$).} Биологически естественно и сделало бы метку
        заметной в данных о стадиях. Отвергнуто, потому что даёт двойной учёт с $B_{\text{std}}$ и
        разрушает точную калибровку распределения стадий.
  \item \textbf{Непрерывная латентная оценка агрессивности.} Реалистичнее --- агрессивность есть спектр, а не
        флаг --- и позволила бы шансам излечения меняться плавно; её нельзя вывести из бинарной колонки, и
        она добавила бы неидентифицируемое латентное измерение.
  \item \textbf{Модель излечения-смеси с меткой в качестве ковариаты.} Статистически стандартная
        формулировка, и текущая конструкция фактически является её логистическим ядром; полная модель-смесь
        оценивала бы долю неизлеченных напрямую, ценой функции правдоподобия, требующей данных наблюдения за
        лечением.
  \item \textbf{Выживаемость по подтипам} (степень злокачественности, статус MSI/MMR, молекулярный подтип)
        из литературы. Существенно информативнее и клинически осмысленнее бинарного флага, и сделала бы
        эффект излечения основанным на доказательствах. Требует данных уровня подтипов, которых в
        агрегированных файлах нет.
  \item \textbf{Влияние агрессивности на риск смерти от рака вместо вероятности излечения.} Эквивалентно по
        действию на смертность, различно по механизму (меняет выживаемость при отсутствии излечения, а не
        шанс излечения). Возможно, более прямо; труднее калибровать и менее интерпретируемо как «успех
        лечения».
  \item \textbf{Отказ от метки.} Проще всего и убирает заданный вручную параметр, но выбрасывает
        единственную информацию о фенотипе опухоли, которую дают входные данные.
\end{itemize}}

\section{Лечение и излечение}

\subsection{Вероятность излечения}

\mech{При клинической диагностике модель вычисляет вероятность излечения
\[
p_{\text{изл}} \;=\; \min\!\Bigl(1,\; \alpha(\text{диаг})\times \theta_{s}\Bigr),
\]
где $\theta_s$ --- эффективность лечения, зависящая от стадии, а $\alpha(\cdot)$ --- зависящий от возраста
множитель эффективности. $\theta = [0.7,\,0.3,\,0.2,\,0.2]$ для стадий I--IV. $\alpha$ --- ступенчатая
функция возраста диагностики с точками разрыва $40/50/60/70$ и значениями
$[1.0,\,0.9,\,0.5,\,0.3,\,0.1]$. Если опухоль агрессивна, \emph{шансы} излечения умножаются на отношение
шансов для данной стадии из раздела~\ref{sec:agg} до обратного пересчёта. Затем излечение разыгрывается
испытанием Бернулли с полученной вероятностью. В эталонной конфигурации $p_{\text{изл}}$ меняется от $0.70$
(стадия I, диагностика в молодом возрасте) до $0.02$ (стадия IV, диагностика после 70 лет).}

\why{Вероятность излечения разложена на стадию и возраст, потому что это две оси, по которым излечимость
действительно наблюдается и публикуется.

\emph{(i) Стадия доминирует, и её измеряют данные.} Стадия при диагностике --- сильнейший одиночный
предиктор выживаемости при раке в любом регистре, и именно по этому измерению организован файл стадий.
Сделать стадию главным фактором --- значит согласовать модель с тем наблюдаемым, которое её ограничивает.

\emph{(ii) Возраст входит ступенчатой функцией сознательно.} Возрастной градиент излечимости реален ---
пожилые хуже переносят лечение, лечатся менее агрессивно и несут большую конкурирующую смертность, --- но
гладкий возрастной эффект потребовал бы параметров формы, которые ничто во входных данных не позволяет
идентифицировать. Четырёхступенчатая функция --- это та гранулярность, в которой публикуется
онкологическая статистика и в которой представляются результаты; она не требует предположений о форме и
прямо понятна клиницисту. Значения ступеней читаются как «относительная эффективность излечения в этой
возрастной группе».

\emph{(iii) Умножение --- правильная композиция.} Оно удерживает результат в $[0,1]$ без ограничения сверху,
делает каждый множитель интерпретируемым как относительный (пациент с наилучшей возрастной эффективностью
получает ровно зависящую от стадии эффективность) и делает отключение эффекта однострочным изменением (все
возрастные константы равны $1$).

\emph{(iv) Ограничение сверху единицей --- требование корректности,} поскольку механизму Бернулли нужна
вероятность; оно безвредно, потому что эталонные значения к нему не приближаются.

\emph{(v) Преобразование через отношение шансов для агрессивной метки ведёт себя разумно.} Оно монотонно и
ограничено, поэтому влияние метки максимально при промежуточных вероятностях излечения и убывает к нулю на
обоих краях, где излечение уже почти несомненно или почти безнадёжно.}

\subsection{Что означает «излечен» для смерти}

\mech{Излеченному пациенту присваивается возраст смерти после излечения, но решение о том, \emph{когда}
человек умирает, принимает правило конкурирующих рисков из раздела~\ref{sec:death}: агент покидает
популяцию в своём естественном возрасте смерти, а причина относится в слое статистики. Поэтому каждый год
модель сравнивает два сохранённых возраста и считает смерть от рака, когда у агента есть рак, он не излечен
(ни клинически, ни при скрининге), возраст смерти от рака без лечения равен текущему возрасту и не
превышает естественный возраст смерти. Возраст смерти от рака у неизлеченного пациента --- это последнее
пересечение стадии плюс экспоненциальная величина для терминальной фазы,
$\delta \sim \mathrm{Exp}(\lambda_{\text{терм}})$, где $\lambda_{\text{терм}}$ подогнано под наблюдаемую
смертность от рака ($0.35$ в год в эталонной подгонке, то есть средняя терминальная выживаемость около
$2.9$ года).}

\why{Определение излечения как «рак больше не определяет ни дату, ни причину смерти» --- именно то, что
сохраняет арифметику скрининга точной и честной.

\emph{(i) Парный контрфактический сценарий остаётся внутри одного агента.} Поскольку возраст смерти от рака
без лечения хранится для каждого пациента, излечение есть просто смена того, какую причину считает
статистика. Выходные показатели числа спасённых жизней и сэкономленных лет оказываются действительно
внутрииндивидуальными разностями, а не разностями между двумя популяциями с разным выборочным шумом, --- а
именно это и делает их осмысленными при используемых размерах выборки.

\emph{(ii) Модель выживаемости после излечения не нужна.} Оставшаяся жизнь излеченного пациента --- это
оставшаяся жизнь популяции, то есть ровно то предположение о конкурирующих рисках, которое уже сделано для
всех остальных. Введение чего-либо иного потребовало бы специфичной для излечения смертности (токсичность
лечения, поздний рецидив) и, следовательно, данных.

\emph{(iii) Сроки смерти от рака несёт терминальная фаза, и это единственная величина, связанная с лечением,
которая действительно подгоняется.} Смертность от рака в данных определяется в первую очередь тем, как
долго живут неизлеченные пациенты после диагностики, поэтому один подогнанный риск для терминальной фазы
--- самый экономичный способ согласовать этот наблюдаемый показатель.}

\why{\emph{Смешение, о котором стоит сказать ещё раз, потому что оно важно.} Поскольку вероятности излечения
заданы вручную, а подгоняется только терминальный риск, наблюдаемая смертность от рака ограничивает
\emph{произведение} двух величин, а не каждую в отдельности. Сценарий с большим числом излечений и более
быстрой смертью от рака может воспроизвести ту же кривую смертности, что и сценарий с меньшим числом
излечений и более медленной смертью. Следствия: доли излечения, о которых сообщает модель, следует читать
как сценарные предположения, а не как оценки; подогнанный терминальный риск поглощает то, в чём ошибаются
принятые вероятности излечения; и только выживаемость по стадиям с наблюдением за пациентами могла бы их
разделить. Это второй плоский гребень в модели, и он заслуживает такой же заметности, как гребень времени до
диагностики из раздела~\ref{sec:nonident}.}

\alt{\begin{itemize}[nosep]
  \item \textbf{Таблицы относительной выживаемости по стадиям, используемые напрямую} (классический подход в
        моделях экономики онкологических программ). Наблюдаемая величина становится входом, механизм
        излечения не нужен, а опубликованная выживаемость воспроизводится точно; но механизм теряется,
        спросить «почему» пациент живёт дольше нельзя, а улучшения лечения можно представить лишь заменой
        таблицы.
  \item \textbf{Модель излечения-смеси с оценкой доли излеченных по данным наблюдения.} Статистически
        стандартная формулировка, и текущая конструкция фактически является её логистическим ядром; требует
        индивидуального наблюдения и позволила бы идентифицировать доли излечения, которые сейчас
        предполагаются.
  \item \textbf{Моделирование избыточной смертности (относительной выживаемости) в виде риска.} Риск,
        относимый на счёт рака, как функция стадии, возраста и времени с момента диагностики вместо
        вероятности излечения. Ближе к тому, как ведут себя данные регистров, и снимает бинарную абстракцию
        «излечен / не излечен»; нужны таблицы дожития населения по возрасту, полу и периоду, а флаг излечения
        --- и вместе с ним учёт спасённых жизней --- становится лишь приближённым.
  \item \textbf{Моделирование маршрутов лечения} (хирургия, химиотерапия, лучевая терапия, каждый со своим
        распределением исходов). Клинически наиболее информативно и естественное место для нескольких
        неактивных параметров (осложнения, смертность от лечения); требует данных о лечении, которых в
        комплекте нет.
  \item \textbf{Поправка на смертность после излечения} (множитель к естественному риску смерти для
        излеченных пациентов, отражающий токсичность лечения или поздний рецидив). Небольшое честное
        расширение; требует данных для калибровки множителя и нарушает простоту «излеченные пациенты ---
        обычные люди».
  \item \textbf{Строгое излечение с полным удалением риска смерти от рака} и повторным розыгрышем смертности
        от прочих причин по таблице дожития с удалённой причиной. Внутренне наиболее чистая формулировка;
        требует таблицы дожития с удалённой причиной, обсуждавшейся в разделе 3.
\end{itemize}}

\section{Скрининг: график, участие, характеристики теста}

\subsection{График приглашений и участие}

\mech{Раунд скрининга проводится в моделируемом году $y$, если $y \ge \id{screening\_date}$ и разность
$(y - \id{screening\_date})$ кратна \id{frequency} (по умолчанию: программа стартует в 10-й моделируемый год и
работает раз в 2 года). В раунде агент \emph{подлежит} обследованию, если он жив, его возраст лежит в
$[\id{start\_age}, \id{finish\_age}]$ (по умолчанию 50--60), он ещё не был обнаружен при скрининге и ещё не
диагностирован клинически (текущий возраст ниже возраста диагностики). Каждый подлежащий агент проходит тест
с вероятностью \id{participation\_rate} (0.5), разыгрываемой независимо в каждом раунде.}

\why{График выражается возрастным окном плюс периодичностью, потому что именно так программы скрининга
задаются и оцениваются, и он точно представим на годовой сетке (раздел 2.2). Два правила подлежания несут
содержательный смысл. Исключение уже обнаруженных агентов упрощает историю скринингового пути агента и
отражает клиническую реальность того, что человек с найденным раком покидает скрининговый маршрут.
Исключение агентов, уже прошедших возраст клинической диагностики, гарантирует, что обнаруженный при
скрининге рак действительно является раком, пойманным \emph{раньше}, --- а это единственный механизм, которым
скрининг способен помочь в этой модели, --- и не позволяет человеку быть учтённым одновременно и как
обнаруженный при скрининге, и как клинически диагностированный.}

\why{Независимое участие в каждом раунде --- предположение с максимальной энтропией при одном агрегированном
показателе посещаемости, и единственное, которое такой показатель способен обосновать. Оно также
воспроизводит реалистичный феномен прерывистого участия: человек может прийти в 52 и пропустить 54. Чего оно
\emph{не} может представить --- это хорошо документированную подгруппу никогда не приходящих (и
коррелированных сверхприходящих). Это подлинное упрощение: польза реальных программ зависит от охвата
недообследованных, и модель с независимым участием завысит популяционный эффект заданного среднего уровня
участия.}

\subsection{Характеристики теста}

\mech{Тест описывается долей истинно положительных и долей ложноположительных результатов, выбираемых парой
из меню по индексу \id{selected\_test}: $\{50,60,70,80\}$ отображаются в
$\text{TP} = \{0.6,0.7,0.9,0.95\}$ и $\text{FP} = \{0.05,0.07,0.1,0.15\}$. Фактически применяется только доля
истинно положительных: агент с доклиническим раком даёт положительный результат с вероятностью TP, а агент
без доклинического рака --- с вероятностью $0.1\times\text{TP}$.}

\why{Меню рабочих точек (TP, FP) делает компромисс по точности явным и исключает физически невозможное
сочетание идеального теста: пользователь может выбрать только из предложенной эффективной границы.
Представление теста через TP/FP, применяемые к двум подгруппам (с доклиническим раком и без него), --- это
классическая формулировка таблицы 2$\times$2, не требующая предположений о распределении значений теста, что
важно, поскольку данных о значениях теста в комплекте нет.}

\why{\emph{Две задокументированные несогласованности в этой области, и обе следует исправить, а не
объяснять.} \emph{(i)} Настроенная доля ложноположительных результатов не работает: вероятность
ложноположительного результата в модели равна $0.1\times\text{TP}$, тогда как \id{test\_fp} читается,
хранится, показывается в интерфейсе и никогда не применяется. В результате смена теста в меню меняет долю
ложноположительных результатов только через корреляцию с TP и всегда при фиксированном отношении 10\%, ---
то есть колонка FP в меню декоративна. \emph{(ii)} Отношение шансов излечения для агрессивной опухоли
(раздел~\ref{sec:agg}) применяется на клиническом пути излечения, но \emph{не} на пути излечения при
скрининге, поэтому агрессивная опухоль, обнаруженная при скрининге, в этом одном отношении трактуется как
неагрессивная. Обе несогласованности малы по величине, и ни одна не ломает модель, но они означают, что
скрининговый модуль не полностью следует собственному набору параметров.}

\subsection{Что делает обнаружение при скрининге}

\mech{Если у подлежащего и участвующего агента есть доклинический рак (рак есть и возраст начала уже
пройден), положительный тест помечает его как обнаруженного при скрининге, фиксирует возраст скрининга и
определяет стадию на момент обнаружения подсчётом того, сколько возрастов пересечения уже пройдено к этому
моменту (тот же механизм вывода стадии, что и при клинической диагностике, раздел~\ref{sec:stage}). Затем
излечение разыгрывается по \emph{той же} модели излечения, что и при клинической диагностике --- возрастной
множитель эффективности умножается на зависящую от стадии эффективность лечения, --- но независимым
розыгрышем. Если розыгрыш даёт излечение, пациент помечается излеченным при скрининге и получает возраст
смерти при скрининге не раньше, чем через пять-десять лет после обнаружения; если излечения нет, флаг
клинического излечения \emph{сбрасывается}, а терминальный возраст без лечения остаётся в силе.}

\why{Важнейшее конструктивное решение здесь --- то, что обнаружение при скрининге переиспользует
клиническую модель излечения с новым розыгрышем, и что вероятность излечения зависит от стадии, \emph{на
которой рак был найден}. Именно отсюда в этой модели берётся польза скрининга: более раннее обнаружение
означает более раннюю стадию, а более ранняя стадия --- более высокую вероятность излечения ($0.7$ на
стадии I против $0.2$ на стадии IV). Никакого отдельного эффекта «скрининг лучше» не предполагается ---
польза целиком является следствием сдвига стадии, что и есть честный способ её моделировать, и способ,
опровержимый по результатам исследований.}

\why{\emph{Перезапись, и почему о ней нужно сказать.} Обнаружение при скрининге \emph{заменяет} клиническую
траекторию: неудачное излечение при скрининге сбрасывает флаг клинического излечения, поэтому пациент,
который был бы излечен клинически, может умереть от рака, будучи обнаруженным при скрининге. Парный учёт
спасённых жизней защищён от двойного счёта --- агент считается спасённым только если он излечен при
скрининге \emph{и} не был излечен клинически, --- но перезапись означает, что модель не рассматривает
«пациента, обнаруженного при скрининге, который и так поправился бы» как нейтральный случай. Поэтому
сообщаемый эффект скрининга есть нетто-результат канала пользы (более раннее обнаружение, более низкая
стадия, больше излечений) и канала вреда (перезапись благоприятного клинического исхода). В эталонном
прогоне при скрининге излечены всего 5 пациентов против 22 обнаруженных, что согласуется с тем, что канал
вреда реален, а не является крайним случаем.}

\subsection{Ложноположительные результаты}

\mech{Участвующий агент без доклинического рака даёт положительный результат с вероятностью
$0.1\times\text{TP}$; событие подсчитывается по годам и формирует сообщаемую кривую ложноположительных
результатов. Никаких других последствий нет: ни дообследования, ни стоимости обследования, ни тревоги, ни
влияния на будущее участие.}

\why{Модель считает то, что способна посчитать. Ложноположительные результаты важны для программы скрининга
главным образом через бремя ненужного дообследования, его стоимость и возможность изменения поведения, ---
а в комплекте входных данных нет информации ни об одном из этих аспектов. Их подсчёт сохраняет феномен
видимым в результатах и делает бездействие настроенной доли (предыдущий подраздел) очевидным для всякого,
кто меняет меню тестов, не выдумывая при этом последствий. Сознательно не моделировать известный эффект
можно защитить; молча моделировать его выдуманными числами --- нет.}

\alt{\begin{itemize}[nosep]
  \item \textbf{Непрерывный по времени процесс скрининга с точными моментами.} Снимает дискретизацию
        возрастов скрининга и позволяет обнаруживать опухоль в момент её обнаруживаемости; требует явного
        порога обнаруживаемости, которого данные не дают, и отказывается от чистой годовой сетки.
  \item \textbf{Коррелированное («залипающее») участие, никогда не приходящие, гетерогенность отклика на
        приглашение.} Хорошо обоснованное и важное усовершенствование: польза скрининга сосредоточена среди
        приходящих, и программа с участием 50\% случайной половины ведёт себя иначе, чем со 100\% участием
        половины населения. Требует индивидуальных данных о посещаемости; текущая конструкция этого выразить
        не может.
  \item \textbf{Чувствительность и специфичность вместо TP/FP.} Математически эквивалентно применяемому
        (чувствительность $=$ TP, специфичность $=1-$FP); текущая формулировка выбрана потому, что модель
        применяет тест к агенту, о котором известно, что у него есть доклинический рак, а это и есть
        определение чувствительности.
  \item \textbf{Зависящая от размера чувствительность теста} (тест улучшается по мере роста опухоли). Именно
        так ведут себя реальные тесты, и это усиливает смещение по длительности; требует связи между
        «размером» по Гомпертцу и обнаруживаемостью и данных о вероятности обнаружения по размеру ---
        естественное дополнение к альтернативе с порогом размера из раздела~\ref{sec:nonident}.
  \item \textbf{Моделирование последствий ложноположительного результата} (процедуры дообследования, их
        стоимость и осложнения, тревога, снижение будущего участия). Правильное долгосрочное направление;
        требует медико-экономических и поведенческих данных, которых нет.
  \item \textbf{Реализм лечения в зависимости от маршрута.} Текущее предположение --- что обнаруженный при
        скрининге и клинически диагностированный рак одной стадии лечатся одинаково. У реальных обнаруженных
        при скрининге опухолей биология часто благоприятнее, поэтому это может занижать пользу; альтернативы
        --- параметр излечения по маршруту или выживаемость по маршруту из исследования.
  \item \textbf{Явный канал интервальных раков} (раки, проявляющиеся между раундами). Здесь он возникает
        эндогенно, поскольку риск диагностики продолжает действовать между раундами; сделать его явным
        означает его измерять, что позволили бы данные исследования.
  \item \textbf{Развёртывание программы во времени} вместо включения на 10-й год. Вопрос сценария и
        отчётности, а не структуры.
\end{itemize}}

\section{Отнесение причины смерти и парное сравнение}
\label{sec:death}

\mech{Модель никогда не \emph{разыгрывает} причину смерти: она её \emph{выводит} каждый год из уже
сохранённых возрастов и признаков. Категории таковы: $0$ --- смерть не зафиксирована (жив), $1$ --- смерть не
от рака (естественная), включая пациента с раком, у которого естественный возраст смерти наступает раньше,
$2$ --- смерть от рака на пути естественной истории без лечения, $3$ --- смерть предотвращена скринингом
(пациент умер бы от рака, но был излечен при скрининге), $4$ --- смерть от других причин после клинического
излечения.

Два ряда смертности формируются \emph{одними и теми же} агентами в \emph{одном} прогоне: ряд без лечения,
использующий возраст смерти от рака из естественной истории каждого пациента, и ряд с учётом скрининга,
использующий эффективный возраст смерти (возраст смерти при скрининге для обнаруженных и излеченных при
скрининге, иначе возраст без лечения). Польза скрининга сообщается как \emph{число спасённых жизней} ---
пациенты, излеченные при скрининге, не излеченные клинически, чья смерть от рака без лечения попала бы внутрь
моделируемого окна (3 человека в эталонном прогоне), --- и \emph{число сэкономленных лет жизни}, накопленная
разность двух возрастов смерти (14,4 года в эталонном прогоне, около 4,8 года на одного спасённого).}

\why{\emph{(i) Вывод причины делает категории взаимоисключающими и исчерпывающими по построению.} Каждый
агент попадает ровно в одну категорию каждый год, поэтому счётчики сходятся, и никакой отдельный учёт не
может с ними разойтись.

\emph{(ii) Оба плеча берутся с одних и тех же агентов, поэтому сравнение парное.} Это самое крупное решение по
статистической эффективности во всей модели. Поскольку траектория без лечения уже сохранена для каждого
пациента, эффект скрининга есть внутрииндивидуальная разность, и её дисперсия на порядки меньше дисперсии
разности двух независимо смоделированных популяций. Польза скрининга в несколько жизней на 20\,000 агентов
вообще измерима именно благодаря этой конструкции; при двух независимых плечах она утонула бы в выборочном
шуме.

\emph{(iii) Отсутствие повторного моделирования означает отсутствие артефактов зерна в контрасте.}
Контрфактический сценарий читается из сохранённого числа, а не получается прогоном контрольного сценария,
поэтому контраст не наследует случайность контрольного прогона и не может быть «подстроен» выбором зерна.

\emph{(iv) Оговорка «не излечен клинически» --- правильный контрфактический критерий.} Агент считается
спасённым только если скрининг его излечил \emph{и} он не был бы излечен в любом случае. Без этой оговорки
каждое раннее обнаружение считалось бы успехом, включая пациентов, которые поправились бы и без него, --- а
это и есть ошибка, из-за которой наивные оценки скрининга завышают пользу.}

\why{\emph{Три честные оговорки.}
\emph{(а)} Агенты покидают популяцию только в своём естественном возрасте смерти, поэтому зафиксированная
смерть от рака агента не удаляет (раздел 3.2): возрастная структура следует таблице дожития по всем причинам,
а показатели по конкретной причине используют знаменатель общей популяции. Это стандартный приём для
конкретной причины при предположении о независимости, и он естественная цель для валидации.
\emph{(б)} Одна из ветвей учёта спасённых проверяет, предшествует ли естественная смерть смерти от рака без
лечения, то есть умирал бы человек от рака \emph{или нет}. В эталонном прогоне она пуста и структурно не
может вносить вклад в число лет (поскольку возраст смерти при скрининге никогда не раньше возраста без
лечения), поэтому она способна лишь завысить счётчик людей. Это скрытая несогласованность, а не текущая
ошибка, и её следует устранить.
\emph{(в)} Ограничение спасённых жизней смертями, которые произошли бы внутри моделируемого окна, делает
сообщаемую величину \emph{внутригоризонтной}, а не пожизненной: человек, чья смерть без лечения пришлась бы
на 45-й год 30-летнего прогона, не учитывается. Это занижает пожизненную пользу и делает заголовочную
величину зависящей от горизонта, --- о чём стоит говорить везде, где она цитируется.}

\alt{\begin{itemize}[nosep]
  \item \textbf{Две независимые популяции (когорта со скринингом и без него).} Проще в реализации и легче для
        объяснения; контраст тогда несёт полную дисперсию Монте-Карло, поэтому для той же точности нужно
        гораздо больше агентов, и ответ становится чувствительным к выбору пары зёрен. Кроме того, он не
        может выразить «две возможные жизни одного и того же человека».
  \item \textbf{Парный прогон с тем же зерном при отключённом скрининге.} Даёт парный контраст без хранения
        возрастов без лечения, ценой удвоения времени работы и зависимости от точного совпадения зёрен.
        Текущая конструкция получает ту же информацию бесплатно.
  \item \textbf{Отнесение через избыточную смертность / относительную выживаемость} (смерти от рака как
        превышение над фоновой таблицей дожития, с атрибутивными долями). Стандарт в анализе регистров и не
        требует моделируемого механизма излечения; отказывается от точного парного учёта, требует фоновую
        смертность с удалённой причиной и превращает «спасённых» в вероятность, а не в счёт.
  \item \textbf{Отнесение по вероятности причинности} (смерть частично относится на счёт каждой причины).
        Ближе к тому, как причины смерти реально кодируются; разрушает целочисленные счётчики, на которых
        построены вычислительные и экономические выходы.
  \item \textbf{Пожизненный учёт с экстраполяцией выживаемости} за пределы горизонта. Даёт подлинно
        сопоставимую величину «число спасённых жизней», используемую в отчётах исследований; требует
        экстраполяции за пределы данных и ставки дисконтирования для лет.
  \item \textbf{Исходы, взвешенные по QALY или DALY,} вместо сырых жизней и лет жизни. Правильная валюта для
        сравнения политик; требует весов полезности, ставки дисконтирования и точки зрения (система
        здравоохранения или общество), --- ничего из чего в комплекте данных нет.
\end{itemize}}

\section{Сообщаемые результаты}

\subsection{Подсчёты, показатели и их знаменатели}

\mech{Все результаты накапливаются как счётчики с индексами \emph{года} и \emph{возраста}: возрастное
распределение, число в группе риска, случаи заболевания (агенты, достигающие возраста начала), диагнозы
(агенты, достигающие возраста диагностики) и смерти от рака в обоих рядах --- без лечения и с учётом
скрининга. Показатель в возрасте $a$ есть суммарное число событий в этом возрасте, делённое на суммарное
число человеко-лет в группе риска в этом возрасте, то есть средний годовой показатель за горизонт.
Сознательно строятся две разные кривые типа «заболеваемость»: \emph{заболеваемость} считает
(ненаблюдаемое) доклиническое начало, а \emph{диагностика} --- (наблюдаемое) клиническое выявление.}

\why{Подсчёт на той же сетке, что и данные, позволяет сравнивать выход модели с опубликованными таблицами,
использованными для калибровки, без перегруппировки, и именно это делает калибровку и проверки вообще
возможными. Разделение двух кривых типа «заболеваемость» --- небольшое конструктивное решение с большой
диагностической отдачей: поскольку кривая начала --- моделируемая величина, а кривая диагностики ---
калиброванная, их расхождение есть просто принятое время до диагностики, сделанное видимым. Поэтому каждый
может увидеть, какая часть выхода модели зависит от ненаблюдаемой величины, --- а в этом и смысл анализа
чувствительности по времени до диагностики.}

\why{\emph{Знаменатель --- общая популяция, а не популяция без рака.} Поскольку агенты удаляются только по
естественной смерти (раздел 3.2), знаменатель группы риска для показателей по конкретной причине включает
людей с зафиксированной смертью от рака. Это стандартная конструкция для конкретной причины при
предположении о независимости; поэтому сообщаемые показатели не следует интерпретировать как «вероятность
смерти от рака среди пациентов с раком». Возрастные показатели приводятся как грубые, поэтому пользователь,
сравнивающий две популяции, должен сначала провести стандартизацию по возрасту.}

\alt{\begin{itemize}[nosep]
  \item \textbf{Стандартизованные по возрасту показатели} (прямая стандартизация к эталонной популяции)
        вместо грубых. Необходимы всегда, когда сравниваются данные с разной возрастной структурой, ---
        а это обычный случай здесь, поскольку наборы данных в комплекте различаются по странам. Выполнение
        стандартизации в модели, а не оставление её читателю, --- небольшое и очень полезное улучшение.
  \item \textbf{Кумулятивный / пожизненный риск} («вероятность диагноза до 75 лет»). Заголовочный показатель
        в большинстве онкологических статистик и более понятный неспециалисту, чем годовой показатель; это
        перевыражение того, что модель уже считает.
  \item \textbf{Ориентация на заболеваемость или на смертность} при подаче результатов. Эпидемиологическая
        конвенция --- формулировать результат либо как предотвращённые случаи, либо как предотвращённые
        смерти; модель сообщает и то, и другое, что является надмножеством, но выбор важен, когда
        последующая модель принятия решений ожидает конкретную валюту.
  \item \textbf{Интервальные, а не одногодовые подсчёты} (пятилетние возрастные группы, календарные
        периоды). Согласуется с тем, как публикуется статистика, и снижает разреженность малых счётчиков;
        теряется разрешение именно там, где модель наиболее информативна (окно скрининга).
  \item \textbf{Сглаживание показателей} (как это делается для кривых в анализе чувствительности). Делает
        графики читаемыми, но меняет числа; оно должно оставаться выбором уровня представления, как сейчас и
        есть.
\end{itemize}}

\subsection{Выживаемость с разбивкой по стадии и агрессивности}
\label{sec:survival}

\mech{Строятся две общие кривые выживаемости по конкретной причине --- одна по ряду без лечения, другая по
ряду с учётом скрининга, --- обе от момента диагностики. Кроме того, модель сообщает выживаемость по
конкретной причине от момента диагностики \emph{с разбивкой по стадии и по агрессивности} (каждая из стадий
I--IV плюс объединённая кривая, в подгруппах агрессивных и неагрессивных опухолей), вместе с \emph{долей
излеченных} в каждой из этих страт. Событие определено точно: пациент умирает от рака, если он не излечен
(клинически или при скрининге) и возраст смерти от рака без лечения наступает раньше естественного возраста
смерти. Излеченные пациенты, а также неизлеченные, умершие раньше от чего-то другого, цензурируются. Кривая
накапливается как дискретный риск, $H_t = H_{t-1} + d_t/n_t$ для числа событий $d_t$ и числа в группе риска
$n_t$, с $S_t = \exp(-H_t)$ на горизонте 30 лет. В эталонном прогоне объединённая доля излеченных равна
$13.1\%$ для агрессивных и $27.2\%$ для неагрессивных пациентов.}

\why{\emph{(i) По конкретной причине, а не общая,} потому что интересующая величина --- это эффект рака, и
потому что общая выживаемость примешивает более крупную смертность не от рака, --- а она различается между
наборами данных в комплекте, так что общая кривая сделала бы межстрановое сравнение бессмысленным.

\emph{(ii) Отсчёт от момента диагностики,} поскольку это клинически осмысленная точка начала и именно она
используется в регистровой отчётности о выживаемости, поэтому выход напрямую сопоставим с опубликованными
величинами.

\emph{(iii) Соглашение «излечен $\Rightarrow$ цензурирован» --- правильный приём для конкретной причины.} Как
только пациент излечен, рак больше не может его убить, поэтому для кривой, специфичной для рака, такой
пациент больше не находится в группе риска события. Считать излеченного пациента событием было бы просто
неверно; трактовать его как цензурированного --- стандартная конструкция, и это то же правило, которым
пользуется слой отнесения причин смерти (раздел 10), поэтому они не могут разойтись.

\emph{(iv) Разбивка по стратам --- защита от парадокса Симпсона.} Поскольку агрессивность резко растёт со
стадией ($0.096$ на стадии I до $0.639$ на стадии IV, раздел~\ref{sec:agg}), объединённые кривые
«агрессивные против неагрессивных» в основном измеряют состав по стадиям. Только сравнение внутри стадии
выделяет предположение об отношении шансов, поэтому публикация одних объединённых кривых приглашала бы
именно к неверному выводу.

\emph{(v) Публикация долей излечения рядом с кривыми делает центральный механизм проверяемым.} Канал пользы
в модели --- «более ранняя стадия $\rightarrow$ выше вероятность излечения», поэтому читатель должен видеть,
скольких пациентов абстракция излечения фактически спасает, по стадии и по агрессивности, --- и проверять,
что сообщаемые доли движутся в направлении, которое подразумевает параметр.}

\why{\emph{Оговорка, вытекающая из раздела 8.} Эти кривые условны по принятым вероятностям излечения,
которые заданы вручную и смешаны с подогнанным терминальным риском. Поэтому кривая выживаемости может быть
«правильной» по частично компенсирующим друг друга причинам. Честное прочтение таково: кривые являются
диагностикой внутреннего механизма модели и инструментом сравнения сценариев, а не независимой оценкой
выживаемости при раке. Сравнение их с внешне опубликованной выживаемостью по стадиям --- то испытание,
которое решило бы вопрос.}

\alt{\begin{itemize}[nosep]
  \item \textbf{Общая выживаемость.} Однозначна --- не требует суждения о причине смерти, --- но смешана с
        фоновой смертностью, поэтому несопоставима между популяциями с разными таблицами дожития.
  \item \textbf{Относительная / чистая выживаемость} (наблюдаемая выживаемость, делённая на ожидаемую по
        фоновой таблице дожития). Регистровый стандарт и единственная конструкция, делающая выживаемость в
        модели сопоставимой с опубликованными величинами; требует фоновой смертности по возрасту, полу и
        периоду.
  \item \textbf{Моделирование выживаемости через избыточный риск} (риск, относимый на счёт рака, как функция
        стадии, возраста и времени с момента диагностики). Ближе к тому, как ведут себя данные, и снимает
        необходимость в бинарном излечении; требует индивидуального наблюдения для оценки.
  \item \textbf{Подогнанная параметрическая модель выживаемости} (Вейбулл, Гомпертц, логнормальная) вместо
        непараметрической кривой. Компактные сводные статистики, экстраполяция и простое сравнение моделей;
        добавляет предположение о форме, которое данные не всегда могут обосновать.
  \item \textbf{Условная выживаемость} (выживаемость при условии, что пять лет уже прожиты). Клинически
        наиболее интерпретирующая формулировка «плато излечения», делающая абстракцию излечения видимой без
        добавления параметра.
  \item \textbf{Отношения рисков между плечами} вместо полных кривых. Компактная мера эффекта, удобная для
        таблиц; теряет временную структуру, а именно её скрининг и затрагивает сильнее всего.
  \item \textbf{Разбивка по дополнительным осям} (возрастная группа, маршрут обнаружения, набор данных).
        Дёшево в вычислении и обнажило бы взаимодействия, которые объединённая отчётность скрывает; риск ---
        малые страты и шумные кривые, поэтому горизонтальное разрешение следует оставить на стадии и
        агрессивности.
\end{itemize}}

\subsection{Сдвиг стадии: почему два распределения по стадиям несопоставимы}
\label{sec:stageshift}

\mech{Модель сообщает распределение по стадиям для клинически диагностированных раков и отдельно для
обнаруженных при скрининге. В эталонном прогоне клиническое распределение равно $0.225/0.250/0.284/0.241$
(средняя стадия $2.55$), тогда как распределение обнаруженных при скрининге ($n = 22$) равно
$0.000/0.318/0.500/0.182$ (средняя стадия $2.86$). Однако для этих же 22 агентов стадия при скрининге
\emph{минус} клиническая стадия равна в среднем $-0.68$: 13 из 22 были обнаружены на более ранней стадии
(трое --- на две стадии раньше, десять --- на одну), восемь --- на той же, один --- позже, а средний интервал
между обнаружением и клиническим диагнозом, который произошёл бы иначе, составляет $4.8$ года.}

\why{Эти числа нужно читать вместе, потому что два распределения по стадиям \emph{несопоставимы} и
публикация только маргинальных распределений провоцировала бы неверный вывод в любую сторону. Скрининг может
поймать только того, кто в момент раунда находится в доклинической фазе, а опухоль с долгим пребыванием имеет
больше возможностей быть пойманной. Это смещение по длительности: группа обнаруженных при скрининге
обогащена медленно растущими опухолями, у которых было \emph{больше} времени на прогрессию. Поэтому вполне
возможно --- и происходит в эталонном прогоне, --- что распределение обнаруженных при скрининге в поперечном
срезе выглядит хуже клинического, тогда как каждое внутрииндивидуальное сравнение показывает, что
обнаруженная при скрининге опухоль найдена раньше, чем была бы найдена иначе.

Механизм пользы в модели --- именно внутрииндивидуальный сдвиг: более низкая стадия при обнаружении означает
более высокую вероятность излечения (раздел~\ref{sec:agg}, $0.7$ на стадии I против $0.2$ на стадии IV).
Поэтому измерение механизма требует сравнения каждого обнаруженного агента с его собственным контрфактическим
сценарием, что модель может сделать, поскольку хранит клиническую стадию для каждого агента. Средний
интервал «обнаружение --- клинический диагноз» --- здесь около $4.8$ года --- есть величина, ограничивающая
достижимую пользу: это тот запас времени, на который скрининговый маршрут действует раньше.}

\why{\emph{Предупреждение о размере выборки.} Двадцать два обнаруженных при скрининге рака на 20\,000
агентов --- малое число, поэтому распределение по стадиям среди обнаруженных несёт широкую
неопределённость и не должно цитироваться как точечная оценка без интервала или без ансамбля по зёрнам.
Средний внутрииндивидуальный сдвиг устойчивее, потому что это парная величина, --- по той же причине, по
которой парная конструкция выбрана для учёта пользы (раздел 10). Это ещё и практический аргумент в пользу
стандартизованного автоматического перебора, описанного в разделе~\ref{sec:sens}: статистика скрининга из
одного прогона слишком шумна, чтобы её просматривать по одному прогону за раз.}

\alt{\begin{itemize}[nosep]
  \item \textbf{Совокупный сдвиг стадии по всем ракам} (обнаруженные при скрининге и остальные вместе). Проще
        всего представить и наименее информативно, потому что это ровно то сравнение, которое искажает
        смещение по длительности.
  \item \textbf{Парный внутрииндивидуальный сдвиг стадии как полноценный выходной показатель} (его среднее,
        медиана и полное распределение, плюс интервал «обнаружение --- диагноз»). Это и есть рекомендация: это
        величина, которая действительно определяет пользу, она устойчива, поскольку парная, и её нельзя
        прочитать как популяционный сдвиг.
  \item \textbf{Распределение времени до диагностики при скрининге} вместо его среднего. Показывает хвост
        очень долгих пребываний, который и определяет совокупное распределение по стадиям.
  \item \textbf{Сдвиг по TNM} (размер опухоли, поражение узлов, отдалённые метастазы) вместо четырёх стадий.
        Клинически осмысленное описание, которое позволило бы сравнивать сдвиг стадии в модели с результатами
        исследований; требует данных о размере и узлах.
  \item \textbf{Стандартизованная по возрасту и стадии заболеваемость и смертность по маршруту обнаружения}
        (классический анализ сдвига стадии, применяемый в литературе о скрининге). Стандарт, который сделал
        бы выход модели сопоставимым с опубликованными оценками скрининга; требует стандартизации и
        сопоставимого знаменателя, а это несложные дополнения.
  \item \textbf{«Доля обнаруженных на стадии I--II» как заголовочный показатель скрининга.} Интуитивно,
        просто для изложения и сразу обнажает оговорку о смещении по длительности, --- что является либо
        преимуществом, либо ловушкой в зависимости от формулировки.
\end{itemize}}

\section{Затраты}

\mech{Модель затрат --- это три строки арифметики. Затраты на скрининг равны числу проведённых тестов,
умноженному на единичную цену теста. Затраты на лечение равны числу \emph{излеченных} пациентов на каждой
стадии, умноженному на единичную цену лечения для этой стадии и просуммированному по стадиям ($1,2,3,4$ по
стадиям в конфигурации из комплекта). Итог --- их сумма. В эталонном прогоне: $13\,111$ тестов, затраты на
скрининг $13\,111$, затраты на лечение $398$, всего $13\,509$.}

\why{Цены безразмерны, потому что достоверных данных о стоимости в комплекте нет, а модели скрининга,
главное сообщение которой --- \emph{ранжирование} стратегий, нужны защитимые относительные, а не абсолютные
цены. Кодирование «тест стоит одну единицу, лечение --- от одной до четырёх единиц в зависимости от стадии»
улавливает только порядковый и бесспорный клинический факт --- лечение на поздней стадии дороже, --- и
оставляет пользователю свободу подставить реальные тарифы. По той же причине калибровка затрат не касается:
они не могут молча поглотить расхождение в другом месте.

Ограничение модуля тремя строками --- само по себе конструктивное решение. Более богатая модель затрат
должна была бы принять предположения (ставка дисконтирования, точка зрения, протоколы наблюдения,
осложнения), которые данные комплекта не могут обосновать, и эти предположения затем определяли бы сравнение
затрат, --- тогда как здесь арифметику можно проверить вручную, а чувствительность ранжирования к ценам
прозрачна.}

\why{\emph{Три ограничения, которые следует читать как оговорки к любой цифре затрат.}
\emph{(а)} Лечение оплачивается только для излеченных пациентов, поэтому с точки зрения плательщика истинные
затраты программы скрининга занижены (пациенты, которых лечили, но которые умерли, не оплачены), а арифметика
«затраты на спасённую жизнь» не является коэффициентом экономической эффективности в обычном смысле.
\emph{(б)} Три из четырёх ценовых параметров не работают: \id{screening\_price} (стоимость программы на
человека), \id{complications\_price} и \id{reoccurrence\_cost} объявлены, хранятся, показываются в
интерфейсе и никогда не используются моделью. Обращение с набором параметров как с конфигурацией, часть
которой декоративна, --- ловушка для пользователя, который меняет цену и не видит никакого эффекта.
\emph{(в)} Нет дисконтирования, нет временных предпочтений, нет заявления о точке зрения и нет стоимости
ложноположительного результата, хотя именно он является самым частым дорогостоящим событием реальной
программы.}

\alt{\begin{itemize}[nosep]
  \item \textbf{Полный анализ «затраты --- эффективность»} (затраты на спасённый год жизни или на QALY,
        дисконтирование, инкрементальные коэффициенты относительно компаратора, пороги готовности платить,
        вероятностный анализ чувствительности по ценам). Это то, чего ожидает потребитель
        медико-экономических результатов, и модель уже вычисляет оба необходимых компонента (затраты и
        сэкономленные годы жизни). Требует точки зрения, ставки дисконтирования, весов полезности и
        стратегии сравнения, --- а это внешние решения, а не данные.
  \item \textbf{Стоимость по маршрутам лечения} (диагностическое обследование, хирургия, химиотерапия,
        лучевая терапия, паллиативная и терминальная помощь, каждая со своей стоимостью по состоянию).
        Стандарт в моделях принятия решений, и он делает затраты на лечение осмысленно зависящими от стадии
        при диагностике; требует данных о маршрутах и тарифах.
  \item \textbf{Точка зрения системы здравоохранения с реальными тарифами} против \textbf{общественной точки
        зрения} с потерей производительности, неформальным уходом и дорогой. Эти два подхода могут
        различаться в разы, и выбор принадлежит пользователю; модель должна просто указывать, какую точку
        зрения она принимает (сейчас --- ни одну).
  \item \textbf{Микроуровневый учёт стоимости скринингового эпизода} (приглашение, набор для теста,
        лабораторная обработка, процедура дообследования при положительном результате). Основной
        реальный драйвер затрат на марже и естественное место для ныне неактивного \id{screening\_price} и
        для непромоделированных последствий ложноположительных результатов.
  \item \textbf{Оплата всех пролеченных, а не только излеченных.} Изменение в одну строку, которое заставило
        бы показатель затрат вести себя так, как ожидает плательщик, и сделало бы программу скрининга
        соответственно менее привлекательной.
  \item \textbf{Стоимость по годам с момента диагностики} (первый год лечения стоит значительно больше
        последующих). Хорошо установленная закономерность в литературе, и она была бы важна для любого
        сравнения, сдвигающего диагнозы во времени.
\end{itemize}}

\section{Элементы, которые существуют, но ничего не делают}
\label{sec:inactive}

\mech{Конфигурация и интерфейсы предоставляют ряд элементов, не влияющих ни на один результат. Они
перечислены здесь, а не тихо опущены, потому что пользователь, меняющий параметр и не видящий никакого
эффекта, введён в заблуждение моделью, а не собственной ошибкой. Каждый пункт ниже проверен поиском
фактических обращений к нему в пути исполнения симуляции.

\begin{center}
\footnotesize\setlength{\tabcolsep}{3pt}
\begin{tabular}{@{}p{4.3cm}p{3.8cm}p{7.2cm}@{}}
\toprule
\textbf{Элемент} & \textbf{Для чего предназначался} & \textbf{Фактический статус} \\
\midrule
Вероятность рецидива & рецидив после излечения & Только разыгрывает экспортируемый флаг
(\id{reoccurrence\_probability} $=0.2$); не влияет на выживаемость, сроки смерти и затраты. \\
Смертность после лечения & дополнительная смертность после лечения & Объявлено и не читается
(\id{treatment\_mortality\_adjustment} $=1.0$). \\
Осложнения лечения & частота осложнений & Объявлено и не читается
(\id{complications} $=[0.1]$). \\
Стоимость осложнения & стоимость на одно осложнение & Объявлено и не читается
(\id{complications\_price} $=1.0$). \\
Стоимость рецидива & стоимость на один рецидив & Объявлено и не читается
(\id{reoccurrence\_cost} $=1.0$). \\
Стоимость программы на человека & стоимость скрининга на человека & Объявлено и не читается
(\id{screening\_price} $=1.0$); оплачивается только цена теста. \\
Доля ложноположительных & доля ложноположительных результатов теста & Читается, хранится и показывается, но
никогда не применяется (\id{test\_fp}, \id{test\_fp\_selected}): в модели используется
$0.1\times$TP. \\
Факторы риска & относительные риски экспозиций & Объявлено и не читается симуляцией
(\id{factors\_rr}, \id{age\_irrelevant\_factors\_rr}, \id{factor\_age\_correction}). \\
Доступ к факторам & вспомогательные методы для этих таблиц & Никогда не вызываются (\id{has\_factor},
\id{get\_factor\_rr}). \\
Устаревшая конфигурация & старый формат «ключ--значение» & Достижима только при указании модели на файл
\id{.txt} (\id{from\_legacy\_txt}, \id{\_csharp\_to\_python\_name}); ни одна конфигурация из комплекта её
не использует. \\
Дублирующий помощник излечения & ступенчатая функция излечения по возрасту & Никогда не вызывается
(\id{get\_age\_cure\_efficiency}); симуляция использует собственную копию той же логики. \\
Перевод рисков в распределение & годовые риски $\rightarrow$ распределение возраста смерти & Никогда не
вызывается (\id{risks\_to\_distribution}); модель использует наблюдаемые смерти напрямую. \\
Помощники вычисления кривой & вычисление подогнанной кривой Гомпертца & Никогда не вызываются
(\id{compute\_output}, \id{compute\_output\_train}, \id{compute\_train\_ages}, \id{gompertz\_v},
\id{compute\_stage\_age}); и подгонка, и симуляция используют времена пересечения в замкнутой форме. \\
Векторизованный розыгрыш риска & векторизованный розыгрыш первого успеха & Никогда не вызывается
(\id{PiecewiseHazard.T\_batch}); генератор популяции делает это встроенно. \\
Помощники скалярной выборки & выборка одного значения & Никогда не вызываются
(\id{Distribution.generate\_single}, \id{RandomState.next\_double}, \id{next\_doubles}). \\
\bottomrule
\end{tabular}
\end{center}}

\why{Принципов, оправдывающих такое перечисление вместо молчаливого удаления, два.
Во-первых, \emph{нереализованный механизм должен быть виден именно как нереализованный}. Несколько из этих
пунктов (смертность от лечения, осложнения, рецидив) соответствуют реальным клиническим эффектам, которые
модель \emph{не} представляет; оставляя их в конфигурации в бездействующем виде, легко навести читателя на
мысль, что они работают, а получившаяся модель опаснее меньшей, которая честно признаёт свои рамки.
Во-вторых, \emph{размер бездействующей поверхности есть мера разрыва между замыслом и реализацией} ---
здесь это довольно концентрированный набор функций, связанных с маршрутами лечения и рецидивом, что служит
полезным сигналом о том, где амбиции модели опередили её входные данные.}

\why{\emph{Слабость, которую это обнажает.} Бездействующие параметры --- ловушка удобства использования.
Пользователь, задавший долю ложноположительных результатов, частоту осложнений или цену скрининга и не
увидевший изменений, не может отличить ошибку от отсутствующей функции, а готовность интерфейса показывать
параметр подразумевает, что он работает. В текущем коде есть встроенные комментарии, помечающие часть из них
как неактивные, и это верная интуиция, но комментарий --- не контракт.}

\alt{\begin{itemize}[nosep]
  \item \textbf{Удалить их.} Наиболее чисто: набор параметров, содержащий только работающие параметры, не
        может ввести в заблуждение. Цена --- потеря задокументированного замысла, а это реальная информация
        для того, кто позже будет реализовывать маршруты лечения.
  \item \textbf{Оставить, но пометить в интерфейсе как неактивные} (сделать неактивными визуально или
        снабдить пометкой «не используется в текущей модели» с указанием причины). Сохраняет замысел и
        устраняет ловушку; это минимальное изменение, решающее собственно проблему.
  \item \textbf{Явно сообщать об ошибке, когда неактивному параметру задают значение, отличное от значения по
        умолчанию} (предупреждение при загрузке). Гарантирует, что никто не будет считать вступившую в силу
        настройку работающей.
  \item \textbf{Реализовать маршруты лечения,} превратив рецидив, осложнения и смертность от лечения в
        работающие механизмы. Правильный долгосрочный ответ и единственный, который делает эти параметры
        осмысленными; требует данных о лечении и осложнениях.
  \item \textbf{Вывести устаревший загрузчик из основного пути} (или удалить его), чтобы остался один формат
        конфигурации. Второй путь разбора --- это вторая поверхность для скрытого расхождения: двум
        загрузчикам пришлось бы согласовывать смысл каждого поля, и ничто это не проверяет.
\end{itemize}}

\section{Калибровка}
\label{sec:calib}

\subsection{Три независимых правдоподобия и всё остальное вручную}

\mech{Когда калибровка запрошена, модель решает три независимые задачи подгонки.

\begin{enumerate}[nosep]
  \item \textbf{Риск диагностики} (6 коэффициентов): правдоподобие Пуассона по наблюдаемым возрастным числам
        случаев, $C(a)\sim\mathrm{Poisson}\bigl(h(a)P(a)\bigr)$, где $h(a)$ --- кусочно-логарифмический риск
        из раздела 4. Минимизируется методом L-BFGS-B при ограничениях
        $\pm5$; сходится за доли секунды.
  \item \textbf{Форма и скорость Гомпертца} (3 параметра $K$, $C$, $B_{\text{pop}}$): метод наименьших
        квадратов по трём временам пересечения, \emph{выведенным из распределения по стадиям} (раздел 6.3),
        плюс штрафы, удерживающие ёмкость среды выше терминальной стадии, скорость положительной и мягко
        регуляризующие ёмкость, чтобы подгонка не ушла в вырожденный угол с огромной ёмкостью и крошечной
        скоростью. В эталонной подгонке невязка порядка $10^{-3}$ --- цели воспроизводятся практически точно.
  \item \textbf{Риск терминальной фазы} (1 параметр $\lambda$): правдоподобие Пуассона по наблюдаемым
        возрастным числам смертей от рака, $D(a)\sim\mathrm{Poisson}\bigl(e^{\lambda}C(a)\bigr)$ при
        $C(a)=I(a)P(a)$. Эталонная подгонка даёт $\lambda=-1.05$, то есть риск $0.35$ в год.
\end{enumerate}

Всё остальное --- настраиваемые значения, а не оценки: среднее время до диагностики, $B_{\text{std}}$,
эффективности лечения по стадиям, возрастные константы излечения, отношение шансов для агрессивности,
график скрининга, участие, точность теста и цены. Обращение распределения по стадиям, дающее времена
пересечения, --- это арифметика, а не оптимизация, и оно выполняется до подгонки Гомпертца.}

\why{\emph{(i) Разделение задачи --- это то, что сохраняет предположения предположениями, а оценки
оценками.} Единая совместная подгонка всех параметров была бы неидентифицируема вдоль двух плоских гребней
из разделов 5.2 и 8.2: оптимизатор мог бы обменять время до диагностики на скорость роста или вероятность
излечения на терминальный риск и улучшить совокупную подгонку, уведя параметры сколь угодно далеко.
Разделение подгонки на три корректно поставленные части означает, что каждый параметр закреплён целью,
которая действительно способна его закрепить, --- а величины, которые не может закрепить ни одна цель,
\emph{не подгоняются вообще}: они настраиваются, а затем варьируются в анализе чувствительности. Это
различие, обеспеченное структурой кода, а не дисциплиной, --- самый важный ограничитель.

\emph{(ii) Каждая целевая функция --- правильное правдоподобие для своих данных.} Числа случаев и смертей
--- это пуассоновские счётчики, а не биномиальные доли, и дискретный риск из раздела 4 является естественной
функцией связи. Использование правильного правдоподобия --- не педантизм: в прежней формулировке
использовалось биномиальное правдоподобие для долей и --- гораздо хуже --- смертность калибровалась по
\emph{риску диагностики}, а не по риску смерти от рака, что делало калибровку смертности бессмысленной, хотя
внешне она выглядела правдоподобно.

\emph{(iii) Малые гладкие целевые функции быстры, воспроизводимы и проверяемы.} Шесть параметров и гладкое
правдоподобие сходятся детерминированно из фиксированной начальной точки; один параметр --- это поиск по
прямой. Ничто не подгоняется моделированием, поэтому калибровка занимает миллисекунды, результаты точно
воспроизводимы, а каждое значение правдоподобия можно напечатать и проверить.

\emph{(iv) Структурные ограничения выражены гладкими штрафами.} Каждое пересечение порога должно быть
достижимо ($\exp K > s$), а скорость положительной; выражение этих требований большими квадратичными
штрафами удерживает оптимизатор вне вырожденной области без решателя с ограничениями и без цикла отбраковки.}

\why{\emph{Три методологические слабости, которые следует назвать.}
\emph{(а)} Правдоподобие смертности является \emph{приближением}: оно рассматривает каждый случай как
подверженный постоянному риску смерти, игнорируя, что умирают только неизлеченные и что смерть наступает
после диагностики с задержкой. Поэтому подогнанное $\lambda$ поглощает то, в чём ошибаются принятые
вероятности излечения, --- а это ровно то смешение из раздела 8.2, увиденное со стороны калибровки.
\emph{(б)} Нет разделения на обучение и проверку: модель калибруется и оценивается на одних и тех же
данных. Это самая слабая методологическая точка во всей конструкции, и её устранение является процедурным, а
не вычислительным.
\emph{(в)} Не всё, что \emph{могло бы} быть подогнано, подгоняется. Калибровка по распределению стадий
специфична для набора данных по построению (времена пересечения выводятся из распределения по стадиям
именно этого набора), поэтому между наборами переносимы только риск диагностики и риск смертности. Поэтому
подлинная внешняя валидация должна была бы отложить одну из этих двух величин или отложить другую
популяцию и принять, что калибровка стадий выводится заново, --- а это ограничивает то, что здесь может
означать «валидация», и об этом следует говорить всюду, где делается утверждение о валидации.}

\alt{\begin{itemize}[nosep]
  \item \textbf{Совместная байесовская калибровка с априорными распределениями и проверками апостериорной
        предсказательной способности} (MCMC или ABC). Статистически правильный способ обращения с плоскими
        гребнями: предположения входят как априорные распределения, а их последствия приходят как
        доверительные интервалы, а не как набор повторных прогонов. Стоит часов вместо миллисекунд, и
        заголовочные числа тогда зависят от априорных распределений, которые сами не выведены из данных.
  \item \textbf{Моделируемое максимальное правдоподобие} (вычисление правдоподобия моделированием, а не
        аналитически). Устраняет приближение в смертности и сделало бы гребень «излечение против
        терминального риска» явным, а не скрытым; целевая функция становится шумной и гораздо дороже, что в
        свою очередь требует множества повторений или стохастического оптимизатора.
  \item \textbf{Приближённое байесовское вычисление или сопоставление с историей.} Работает без
        правдоподобия и хорошо подходит симулятору, чьё аналитическое правдоподобие лишь приблизительно;
        требует тысяч прогонов и явного допуска, то есть это исследовательский рабочий процесс, а не
        интерактивный.
  \item \textbf{Суррогатная модель плюс глобальная оптимизация} (пути дифференциальной эволюции и Optuna уже
        присутствуют в коде). Полезны, когда целевая функция многоэкстремальна или шумна; добавляют в задачу
        гиперпараметр --- собственную подгонку суррогата.
  \item \textbf{Ручная экспертная (дельфийская) калибровка,} как это практикуется для ряда крупных
        онкологических моделей. Прозрачна и способна учесть суждение, которого нет в данных; не
        воспроизводима и трудна для проверки.
  \item \textbf{Итеративная последовательная калибровка до сходимости} (подогнать риски, зафиксировать их,
        подогнать стадии, вернуться к рискам). Делает связь между компонентами явной; требует критерия
        сходимости и может дрейфовать, если компоненты связаны слабо.
  \item \textbf{Полноценный протокол отложенной выборки:} калибровка на одной популяции или периоде и
        валидация на другом --- наличие наборов данных по США и Германии в комплекте делает это
        однострочным экспериментом для переносимых компонентов, --- временное разделение там, где есть
        несколько периодов, и публикация ошибки валидации как заголовочного показателя качества. Это самое
        ценное методологическое дополнение, доступное без новых данных.
\end{itemize}}

\section{Случайность и воспроизводимость}

\subsection{Два уровня случайности}

\mech{Модель использует два уровня. \emph{Популяционные} розыгрыши (начальные возрасты, естественные
возрасты смерти, зёрна потоков агентов) берутся из одного генератора PCG64, инициализируемого раз за прогон и
сохраняемого вместе с результатами. \emph{Агентные} розыгрыши внутри ядра истории рака этот генератор вообще
не используют: у каждого агента есть 64-битный идентификатор потока, и каждая величина разыгрывается как
$\id{uniform}(\text{seed}_i + k)$ для последовательных малых смещений $k$ с помощью лавинного хеша в стиле
SplitMix64. Гауссовы величины получаются методом Бокса --- Мюллера по двум равномерным из того же потока;
бернуллиевские --- сравнением равномерной с вероятностью. Состояние генератора является \emph{локальным для
потока}, поэтому параллельные прогоны в одном процессе не могут переинициализировать друг друга, а зерно
пишется в журнал при каждом прогоне.}

\why{\emph{(i) Потоки на агента развязывают агентов между собой и от порядка работы.} Розыгрыши агента
зависят только от его собственного идентификатора, поэтому результаты инвариантны к тому, как разложен цикл,
к векторизации и к любому будущему распараллеливанию. Именно это позволяет ядру перебирать агентов в любом
удобном порядке и компилироваться в плотный цикл.

\emph{(ii) Соседние индексы должны декоррелироваться, потому что ядро читает их подряд.} Каждый агент
разыгрывает много величин по близким смещениям, поэтому отображение индекса в значение должно выглядеть как
независимый шум для смещений 1, 2, 3\ldots. Оно им не было. Прежний линейный конгруэнтный генератор с
единственным шагом xorshift давал $\mathrm{corr}\bigl(u(s),u(s{+}1)\bigr)\approx+0.36$ и
$P\bigl(u(s{+}1)<u(s)\bigr)\approx0.12$ --- вместо $0$ и $0.5$, --- что молча связало розыгрыш стадии со
временем до диагностики, розыгрышем излечения, возрастом смерти от рака и меткой агрессивности. Агрессивные
опухоли, например, разыгрывались почти исключительно на стадии IV. Это была настоящая ошибка модели, а не
эстетическая: она меняла совместное распределение собственных переменных модели. Лавинный хеш это
исправляет, и тест-страж теперь фиксирует инвариант ($|\mathrm{corr}|<0.015$ и $|P-\tfrac12|<0.015$ на
десяти смещениях).

\emph{(iii) Локальное для потока состояние --- требование корректности для веб-сервиса.} Приложение
обслуживает параллельные прогоны из разных потоков; общепроцессный синглтон позволял одному прогону
переинициализировать поток, из которого черпал другой, что ломало воспроизводимость лишь периодически и
поэтому выглядело как нестабильность, а не как ошибка.}

\why{\emph{(iv) PCG64 для популяционных розыгрышей} --- современный генератор с небольшим состоянием и
хорошей проверенной статистикой, а также текущий генератор по умолчанию в NumPy, --- безопасный выбор для
величин, разыгрываемых один раз на агента массово.

\emph{(v) Запись зерна в журнал делает любой результат воспроизводимым одной командой,} что является
предпосылкой для рабочих процессов калибровки, отладки и валидации и для воспроизведения цифры, приведённой
в отчёте.}

\why{\emph{Две честные оговорки.} Во-первых, агентные потоки \emph{псевдо}независимы по построению: они не
являются формально независимыми, и их пригодность опирается на проверенное свойство лавинности (что и делает
тест-страж), а не на теорему. Во-вторых, точное воспроизведение прогона требует согласованной инициализации
обоих уровней, поэтому зерно, сообщаемое вместе с результатом, следует понимать как управляющее и агентными
зёрнами тоже, --- так оно и есть, поскольку они сами разыгрываются из инициализированного популяционного
генератора, но эту зависимость стоит обозначить.}

\alt{\begin{itemize}[nosep]
  \item \textbf{Единый последовательный поток случайных чисел для всего.} Простейшая возможная конструкция;
        она делает каждый результат зависящим от точного порядка операций, что исключает векторизацию и
        распараллеливание и означает, что любая перестановка кода меняет все выходные данные. Она же делает
        случайные корреляции между переменными лёгкими в создании и трудными для обнаружения.
  \item \textbf{Полноценный объект генератора на каждого агента} (со своим состоянием PCG64). Статистически
        превосходно и тривиально независимо; создание объекта на агента стоит на порядки больше, чем
        хеширование двух 64-битных слов, а ядро выполняется для каждого агента популяции.
  \item \textbf{Счётчиковые генераторы (Philox, Threefry).} Созданы именно для такого шаблона --- счётчик
        плюс ключ выдаёт поток значений, --- с документированным статистическим качеством и дешёвыми
        параллельными переходами; естественный путь развития, который убрал бы зависимость от лавинности
        хеша.
  \item \textbf{Общие случайные числа между сценариями} (совместное использование случайных толчков между
        плечом со скринингом и без него для дальнейшего уменьшения дисперсии разности). Парная конструкция
        из раздела 10 уже даёт большую часть этой выгоды, поскольку контрфактический сценарий хранится, а не
        моделируется повторно.
  \item \textbf{Повторные прогоны с сообщаемыми интервалами} вместо одного прогона (усреднение по зёрнам).
        Честный способ представления выхода стохастической модели, особенно для хвостовых величин вроде
        числа спасённых жизней, где счёт одного прогона может сместиться на несколько единиц; умножает
        время работы на число повторений, чего интерактивный интерфейс сейчас позволить себе не может.
  \item \textbf{Антитетическая или квазислучайная (Соболь) выборка} для популяционных розыгрышей. Снижает
        дисперсию для величин типа интегралов; гораздо менее полезна для редких хвостовых событий, а именно
        они и являются интересными результатами.
  \item \textbf{Записываемый журнал розыгрышей} для битовой воспроизводимости при любых будущих изменениях
        кода. Сегодня это избыточно, но именно таков стандарт, которому формально сертифицированная модель
        в конечном счёте должна соответствовать.
\end{itemize}}

\section{Анализ чувствительности как утверждение о неопределённости}
\label{sec:sens}

\mech{Анализ чувствительности перезапускает всю модель для множителей среднего времени до диагностики ---
по умолчанию $0.5$, $0.75$, $1.0$, $1.25$ и $1.5$, --- и показывает кривые заболеваемости, смертности,
смертности с учётом скрининга, выживаемости и сэкономленных лет, а также распределения по стадиям для каждого
множителя рядом друг с другом. Для множителя $f$ среднее время до диагностики становится $f\mu$ \emph{и}
времена пересечения стадий масштабируются тем же множителем, после чего Гомпертц подгоняется заново под
новые цели, так что модель по-прежнему воспроизводит то же распределение стадий при новом масштабе времени.
Все остальные параметры остаются на своих подогнанных эталонных значениях, а случайное зерно одинаково для
всех множителей. Сглаживание применяется к кривым при отображении и никогда к сохранённым числам.}

\why{\emph{(i) Честный ответ на неидентифицируемый параметр --- это диапазон, а не значение.} Время до
диагностики нельзя оценить по доступным данным (раздел~\ref{sec:nonident}), поэтому единственное защитимое
действие с ним --- показать, как выглядит ответ по всему правдоподобному диапазону. Анализ чувствительности
\emph{и есть} утверждение модели о неопределённости её единственного структурного предположения, а не
необязательная надстройка.

\emph{(ii) Масштабирование времён пересечения вместе со временем до диагностики сохраняет сравнение
осмысленным.} Поскольку $t_s=-\mu\ln P(\text{стадия}\ge s)$, изменение $\mu$ при фиксированных временах
пересечения изменило бы распределение стадий и тем самым смешало бы вопрос о масштабе времени с изменением
стадирования. Совместное движение обеих величин удерживает наблюдаемое распределение стадий постоянным,
поэтому кривые различаются \emph{из-за} предположения о масштабе времени и ни из-за чего больше.

\emph{(iii) Общее зерно для всех множителей --- приём снижения дисперсии, и здесь он необходим.} Совместное
использование случайных толчков означает, что различия между множителями есть подлинная реакция модели на
предположение, а шум Монте-Карло в значительной мере сокращается в сравнении. Без этого, и при столь малом
числе событий скрининга, как в эталонном прогоне, разброс между кривыми мог бы определяться выборочной
вариацией: пользователь читал бы шум как чувствительность.

\emph{(iv) Сглаживание существует только для отображения,} поэтому сообщаемые и экспортируемые числа
остаются точными, а путь параллельного исполнения нужен, чтобы пять полных повторных прогонов оставались
интерактивными, с автоматическим порогом, избавляющим от накладных расходов на запуск процессов для слишком
малых задач.}

\why{\emph{Две ограничения, о которых следует знать читателю этих графиков.} Во-первых, \emph{варьируется
только время до диагностики}. Остальные заданные вручную величины --- $B_{\text{std}}$, вероятности
излечения, отношение шансов для агрессивности, участие, точность теста --- столь же непроверены, и набор
кривых, меняющийся лишь по одной оси, легко прочитать как «это единственное слабое предположение».
Во-вторых, анализ сообщает \emph{кривые}, а не показатель для принятия решений, поэтому величина, которую
принимающий решения действительно хотел бы видеть в диапазоне (число спасённых жизней или затраты на
спасённую жизнь при каждом множителе), не сведена в таблицу и должна собираться из блока метрик вручную.}

\alt{\begin{itemize}[nosep]
  \item \textbf{Стандартизованный перебор по всем «мягким» параметрам} с представлением результатов в виде
        таблицы показателей для решений, а не семейств кривых. Анализ по одному фактору даёт нижнюю границу
        неопределённости (он не может показать взаимодействия); его легко запускать, поскольку модель быстра,
        и это самое ценное единичное дополнение.
  \item \textbf{Глобальный анализ чувствительности} (разложение дисперсии, индексы Соболя, отсеивающие
        методы по элементарным эффектам). Распределяет дисперсию выхода по предположениям и обнажает
        взаимодействия; требует спланированного эксперимента на сотни или тысячи прогонов, что для текущей
        модели достижимо, а для текущего интерфейса --- нет.
  \item \textbf{Вероятностный анализ чувствительности} (распределения «мягких» параметров, Монте-Карло по
        розыгрышам, интервалы для исходов). Медико-экономический стандарт, превращающий «вот пять кривых» в
        «вот распределение ответа»; требует защитимых распределений, а для заданных вручную параметров это
        означает априорные распределения.
  \item \textbf{Сценарный анализ} («оптимистичный тест», «пессимистичное участие», «медленные опухоли»)
        вместо числового перебора. Более понятен неспециалистам и заставляет называть предположения по
        имени; это изложение той же информации, теряющее систематический охват.
  \item \textbf{Анализ ценности информации} (какое предположение, будучи исследованным, сильнее всего
        уменьшит неопределённость решения). Естественный следующий вопрос после анализа чувствительности и
        тот, который превращает это упражнение в исследовательскую программу; требует правила принятия
        решений и функции потерь.
  \item \textbf{Структурная чувствительность:} замена самого механизма (модель обнаружения по порогу
        размера вместо модели времени до диагностики, длительности на стадию вместо времён пересечения,
        естественная история, управляемая началом заболевания, вместо управляемой диагностикой) и сравнение
        результатов. Никакой перебор параметров не достигает этого класса неопределённости, а именно в нём
        лежит наибольший неоценённый риск.
\end{itemize}}

\section{Предположения, направление смещения и опровержимость}
\label{sec:assumptions}

\subsection{Реестр предположений}

\mech{Ниже перечислены все упрощения модели с указанием направления, в котором каждое из них смещает
результаты. Там, где направление указано, оно указывается относительно «истинной» модели того же вопроса;
там, где его нельзя определить из одной конструкции, об этом и говорится вместо догадки.

\begin{center}
\footnotesize
\begin{tabular}{@{}p{8.0cm}p{7.6cm}@{}}
\toprule
\textbf{Предположение} & \textbf{Направление смещения, если оно неверно} \\
\midrule
Нет стадии предрака (нет последовательности «аденома--карцинома»; все раки возникают инвазивными). &
Завышает пользу скрининга: удаление предракового образования --- дополнительный защитный механизм, которого в
модели нет. \\
Нет непрогрессирующих образований. & Завышает пользу: каждый обнаруженный при скрининге рак считается таким,
который проявился бы и мог убить, поэтому гипердиагностика непредставима. \\
Естественная история управляется диагностикой (нет смертей от недиагностированного рака). & Занижает
онкологическое бремя популяции; безвредно для контраста по скринингу, вводит в заблуждение в оценке бремени. \\
Абсолютный масштаб доклинического времени (среднее время до диагностики) принят, а не оценён. & Неопределимо
--- именно поэтому он варьируется (раздел~\ref{sec:sens}). Это крупнейшее единичное предположение модели. \\
Неоднородность роста ($B_{\text{std}}$) задана вручную. & Неопределимо; слишком малое делает сдвиг стадий
слишком «чистым» и слишком благоприятным, слишком большое даёт обратное. Сейчас не варьируется. \\
Вероятности излечения заданы вручную и смешаны с подогнанным терминальным риском. & Совокупная смертность
от рака согласуется в любом случае; неидентифицированным остаётся состав --- скольких пациентов абстракция
«излечения» фактически спасает. \\
Агрессивность не влияет на рост. & Занижает биологию; наблюдаемый градиент выживаемости по стадиям
оказывается частично конструкцией модели, а не механизмом. \\
Один первичный рак на человека; нет вторых первичных и метахронных опухолей. & Слегка занижает
заболеваемость и смертность. \\
Нет смертности от лечения и осложнений. & Занижает вредную сторону как лечения, так и скрининга
(дообследование после положительного результата). \\
Нет поведенческих эффектов (ложноположительный результат не меняет посещаемость; скрининг не меняет образ
жизни). & Неопределимо; в реальных программах наблюдаются оба эффекта, и в противоположных направлениях. \\
Одинаковое лечение при данной стадии независимо от сопутствующих заболеваний, доступа и эпохи. & Занижает
вариацию и неравенство; среднее может быть смещено в любую сторону в зависимости от популяции. \\
Замкнутая популяция: нет рождений, нет миграции. & Существенно лишь на горизонтах, достаточных для дрейфа
возрастной структуры, а для 30 лет это в основном не так. \\
Дискретное годовое время; порядок событий внутри года игнорируется. & Небольшое, но систематическое для
самых быстрых терминальных фаз (средняя терминальная выживаемость --- лишь около трёх лет). \\
Смерть от рака не удаляет агента из моделируемой популяции; показатели по конкретной причине используют
знаменатель общей популяции. & Рассматривает две причины как независимые; при таких показателях это
небольшой эффект, но он занижает поправку, которую дал бы знаменатель конкурирующих рисков. \\
Возраст --- единственное измерение экспозиции: нет пола, региона, временного тренда, факторов риска. &
Модель представляет одну популяцию в одном периоде; применять её к трендам, к другому полу или к
риск-стратифицированной программе было бы неверно, а не просто неточно. \\
Польза измеряется внутри горизонта моделирования. & Занижает пожизненную пользу и делает заголовочную
величину зависящей от горизонта. \\
Характеристики теста --- постоянная пара (истинно-, ложноположительные). & Упускает зависимость от размера и
весь вред ложноположительных результатов, поэтому чистая польза слегка завышена. \\
Обнаруженные при скрининге и клинически диагностированные раки одной стадии лечатся одинаково. &
Неоднозначно: обнаруженные при скрининге опухоли часто имеют лучшую биологию (занижает пользу), тогда как
перезапись клинического пути излечения может действовать в обратную сторону (раздел 9.3). \\
Нет дисконтирования и временных предпочтений. & Завышает долгосрочную пользу относительно затрат по
сравнению с дисконтированным анализом. \\
Факторы риска, осложнения, рецидив и смертность от лечения существуют как параметры, но ничего не делают
(раздел~\ref{sec:inactive}). & Влияния нет; риск чисто в том, что читатель сочтёт их работающими. \\
\bottomrule
\end{tabular}
\end{center}}

\subsection{Что показало бы, что модель неверна}

\mech{Доступны три класса проверок, и они сильно различаются по силе.

\emph{(i) Проверки внутренней согласованности} --- величины, которые модель должна воспроизводить по
данным, ей предоставленным: возрастная кривая заболеваемости (подгоняется), возрастная кривая смертности от
рака (подгоняется, но со смешением «излечение против терминального риска» из раздела 8.2) и распределение по
стадиям при диагностике. Информативна третья, поскольку из данных берутся только три её \emph{цели} ---
времена пересечения: само распределение по стадиям затем порождается взаимодействием модели роста и времени
до диагностики. В эталонном прогоне оно воспроизводится с точностью до нескольких тысячных
($0.225/0.250/0.284/0.241$ против $0.220/0.256/0.308/0.218$). Неудача здесь указала бы прямо на модель
роста, обращение времён пересечения или агентные розыгрыши.

\emph{(ii) Внешние проверки} --- те, что действительно способны опровергнуть модель: исследование скрининга
с тем же тестом в том же возрастном диапазоне, где польза в модели должна совпасть и по знаку, и по
величине; наблюдаемый состав и частота интервальных раков в реальной программе; внутрииндивидуальный сдвиг
стадии, измеренный в обследованной когорте (раздел~\ref{sec:stageshift}); и время пребывания, выведенное из
сравнения обнаруженных при скрининге и интервальных раков, что позволило бы идентифицировать единственную
величину, которую модель вынуждена предполагать. Ни одна из них невозможна на агрегированных данных из
комплекта, и все они возможны при наличии данных исследования или регистра.

\emph{(iii) Структурные проверки} --- согласованность механизма с тем, что известно о заболевании. В модели
нет стадии предрака и нет непрогрессирующих образований, поэтому если она предсказывает пользу там, где
исследование её не нашло, эти два упущения --- первые подозреваемые, именно в этом порядке. Структурную
проверку можно провести и как эксперимент, а не как рассуждение: реализовать пропущенный механизм и
посмотреть, устоит ли вывод.}

\why{Опровержимость нельзя приделать в конце: её нужно задать до того, как увидены числа, иначе она
вырождается в постфактумную рационализацию. Модель исключительно хорошо подходит для этого, поскольку она
быстра и детерминирована при заданном зерне --- критерий можно проверить за секунды, --- однако сейчас у неё
нет ни предварительно зарегистрированных критериев успеха, ни формальной валидации против чего-либо вне её
собственных калибровочных данных. Это и есть разрыв между работающей моделью и защитимой.}

\alt{\begin{itemize}[nosep]
  \item \textbf{Предварительно зарегистрированные критерии успеха и план анализа} (какие выходные показатели,
        какой допуск, какой набор данных --- решено заранее). Дёшево, чисто процедурно и превращает
        валидацию из повествования в проверку.
  \item \textbf{Постоянный набор эталонных тестов} (зафиксированные выходные данные при фиксированных
        зёрнах, сверяемые при каждом изменении). Для генератора случайных чисел модель это уже делает;
        распространение той же идеи на всю симуляцию ловило бы непреднамеренные изменения, --- те, что дают
        правдоподобно выглядящие, но другие числа, а они и самые опасные.
  \item \textbf{Периодическая перекалибровка по вновь публикуемым данным} со сравнением предсказаний прежней
        версии с новыми наблюдениями и публикацией расхождения. Превращает ежегодный выпуск данных в
        опубликованный ретротест.
  \item \textbf{Независимая реализация одного из компонентов} (например, независимый пересчёт обращения
        времён пересечения и подгонки Гомпертца), чтобы поймать общее заблуждение, разделяемое спецификацией
        и кодом.
  \item \textbf{Проверка на уровне принятия решений} (меняет ли модель решение, которое иначе было бы
        принято?). Самая релевантная проверка для модели скрининга и единственная, которая прямо отвечает на
        вопрос, важна ли ошибка модели.
\end{itemize}}

\section{Сводка: каждый элемент, его решение и альтернативы}

Таблица ниже --- весь этот документ в одном месте. Для каждого элемента моделирования указаны выбранная
конструкция, рассмотренные альтернативы и то, что изменилось бы выше по цепочке при смене конструкции.

{\footnotesize\setlength{\tabcolsep}{4pt}
\begin{longtable}{@{}p{2.3cm}p{4.5cm}p{4.0cm}p{4.0cm}@{}}
\toprule
\textbf{Элемент} & \textbf{Выбранная конструкция} & \textbf{Основные альтернативы} & \textbf{Что изменится при смене} \\
\midrule
\endfirsthead
\toprule
\textbf{Элемент} & \textbf{Выбранная конструкция} & \textbf{Основные альтернативы} & \textbf{Что изменится при смене} \\
\midrule
\endhead
\bottomrule
\endlastfoot
Класс модели & Микросимуляция на агентах; оба плеча из одних и тех же агентов & Когортная марковская модель,
возрастно-структурированное УрЧП, уравнение восстановления, полумарковская, компартментные ОДУ & Теряется
парный учёт спасённых жизней (статистическая эффективность модели); неоднородность приходится
дискретизировать в состояния. \\
Время & Годовые дискретные шаги на целочисленной сетке возрастов & Непрерывное время с экспоненциальными
часами, месячные шаги, событийное моделирование & Моменты скрининга и сроки терминальной фазы становятся
точными; цели калибровки остаются годовыми, поэтому дополнительное разрешение данными не контролируется. \\
Рамки популяции & Замкнутая поперечная когорта, горизонт 30 лет & Когорта рождения с возраста 0, открытая
популяция с демографическим прогнозом, стационарная & Когорте рождения нужен прогрев длиной в столетие;
открытая популяция загрязняет контраст по скринингу демографическим дрейфом. \\
Начальные возрасты & Наблюдаемая возрастная структура как распределение вероятностей & Равномерный или
параметрический профиль, возрастные группы, микроданные переписи & Параметрические профили навязывают форму;
равномерный старт разрушает главный драйвер заболеваемости. \\
Смертность не от рака & Один разыгранный естественный возраст смерти на агента, по наблюдаемым смертям &
Ежегодные розыгрыши Бернулли, параметрический закон смертности, таблицы с удалённой причиной, хрупкость &
Таблицы с удалённой причиной правильно разделили бы две причины (нужен столбец смертей по причинам);
хрупкость требует неидентифицируемого распределения. \\
Риск начала & Кусочно-постоянный логарифмический риск по возрасту; первый успех Бернулли задаёт возраст
диагностики & Гладкий параметрический риск, сплайны, наблюдаемые показатели напрямую, механистическая модель
канцерогенеза & Сплайны --- естественное развитие; механистическая модель сделала бы факторы риска
моделируемыми, но неидентифицируема по одной заболеваемости. \\
Что считается раком & Диагноз до естественной смерти; агенты старше 98 исключены & Естественная история,
управляемая началом; доля непрогрессирующих образований; никаких исключений & Управление началом сделало бы
гипердиагностику представимой и является самым крупным структурным улучшением. \\
Время до диагностики & Экспоненциальное с принятым средним (3,5 года) & Гамма/Вейбулл, детерминированное,
длительности на стадию, обнаруживаемость по порогу размера, оценка по данным исследования & Масштаб
идентифицируют только данные исследования или регистра; обнаруживаемость по размеру объединила бы
обнаружение с кривой роста. \\
Идентификация масштаба времени & Масштаб фиксируется; времена пересечения выводятся из распределения
стадий; масштаб варьируется & Априорное и апостериорное байесовское распределение, профильное
правдоподобие, фиксация скорости роста, время пребывания с двумя параметрами & Апостериорное распределение
превращает одно число плюс оговорку в интервалы для всех выходов, ценой часов и априорного распределения. \\
Кривая роста & Редуцированный Гомпертц, $V(t)=\exp(K-\exp(C-Bt))$ & Экспонента, логиста, Берталанфи,
Вейбулл/степенная, механистические ОДУ, ветвящийся процесс, без модели роста & Другая кривая меняет времена
пересечения в замкнутой форме, а вместе с ними всю калибровку стадий и аргумент о неидентифицируемости. \\
Неоднородность роста & Логнормальная $B_{\text{ind}}$ при заданном вручную $B_{\text{std}}=0.1$ & Без
неоднородности, гамма-распределённая скорость, дискретные классы роста, случайные эффекты по ковариатам &
$B_{\text{std}}=0$ убирает смещение по длительности и делает стадию детерминированной функцией времени до
диагностики. \\
Стадирование & Выводится: число порогов размера, пересечённых до диагностики & Разыгрывание из наблюдаемого
распределения, порядковая регрессия по ковариатам, латентная многосостоянийная прогрессия, TNM & Разыгрывание
распределения делает калибровку тавтологией; TNM требовал бы данных о размере и узлах. \\
Агрессивность & Бернуллиевская метка по стадиям, меняющая шансы излечения (OR 0.667) & Модификатор роста,
непрерывная латентная оценка, модель излечения-смеси, выживаемость по подтипам, без метки & Как модификатор
роста она даёт двойной учёт с $B_{\text{std}}$ и ломает точную подгонку распределения стадий. \\
Излечение & Эффективность по стадии $\times$ ступенчатая функция возраста, с ограничением, Бернулли; излечение
означает, что рак больше не определяет смерть & Таблицы относительной выживаемости напрямую, модель-смесь по
данным наблюдения, моделирование избыточного риска, маршруты лечения & Таблицы регистров убрали бы механизм и
возможность спросить, почему пациент выживает. \\
График скрининга & Возрастное окно плюс фиксированная периодичность; независимое участие в каждом раунде &
Непрерывный по времени скрининг, коррелированное участие, никогда не приходящие, развёртывание программы &
Коррелированное участие сделало бы пользу зависящей от охвата недообследованных; текущая конструкция этого
выразить не может. \\
Характеристики теста & Меню (TP, FP), из которого применяется только TP (FP $=0.1\times$TP) &
Чувствительность/специфичность, зависящая от размера чувствительность, явные затраты на дообследование &
Оживление FP --- это исправление ошибки; зависящая от размера чувствительность делает смещение по
длительности количественным. \\
Обнаружение при скрининге & Использует ту же модель излечения со стадией на момент обнаружения; перезаписывает
клинический путь & Аддитивный эффект, параметр излечения по маршруту, выживаемость по маршруту из
исследования & Перезапись создаёт канал вреда; аддитивная конструкция убрала бы его, но занизила бы пользу. \\
Ложноположительные & Только подсчитываются, без последствий & Процедуры дообследования, затраты, тревога,
влияние на посещаемость & Их моделирование --- правильное направление, и оно требует медико-экономических и
поведенческих данных. \\
Атрибуция смерти & Выводится из сохранённых возрастов; пять исчерпывающих категорий & Сравнение двух
популяций, парный прогон с тем же зерном, относительная выживаемость, вероятность причинности & Сравнение
двух популяций теряет парность и требует гораздо больше агентов для той же точности. \\
Показатели и исходы & Счётчики по году и возрасту; грубые показатели на человеко-годы & Стандартизация по
возрасту, кумулятивный риск, интервальные группы, сглаживание & Стандартизация необходима для межнаборного
сравнения, и сейчас она оставлена читателю. \\
Выживаемость & По конкретной причине от момента диагностики, с разбивкой по стадии и агрессивности,
излеченные цензурируются & Общая выживаемость, относительная/чистая выживаемость, избыточный риск,
параметрические модели, условная выживаемость & Относительная выживаемость --- регистровый стандарт и
единственное сопоставимое сравнение с опубликованными величинами. \\
Сдвиг стадии & Публикуются оба распределения; механизмом считается внутрииндивидуальный сдвиг & Только
совокупный сдвиг, парный сдвиг как полноценный показатель, сдвиг по TNM, стандартизованный анализ сдвига &
Публикация парного сдвига предотвратила бы неверное прочтение маргинальных распределений из-за смещения по
длительности. \\
Затраты & Тесты $\times$ единичная цена плюс излеченные $\times$ цена по стадии; безразмерно & Полный
анализ «затраты --- эффективность», стоимость по маршрутам лечения, реальные тарифы, общественная точка
зрения, микроуровневый учёт & Оплата всех пролеченных и добавление затрат на дообследование изменили бы
экономику качественно, а не только по уровню. \\
Неактивные элементы & Документируются, а не удаляются (раздел~\ref{sec:inactive}) & Удалить, пометить в
интерфейсе как неактивные, предупреждать при нестандартных значениях, реализовать маршруты лечения &
Реализация рецидива, осложнений и смертности от лечения сделала бы эти параметры осмысленными и добавила бы
реальный канал вреда. \\
Случайность & PCG64 на популяционном уровне, хеш SplitMix64 на агента, локальное для потока состояние &
Единый последовательный поток, генератор на агента, счётчиковые генераторы, общие случайные числа, повторные
прогоны с интервалами & Повторные прогоны с интервалами --- честное представление хвостовых величин вроде
числа спасённых жизней. \\
Калибровка & Три независимых правдоподобия; всё неидентифицируемое настраивается вручную & Совместная
байесовская калибровка, моделируемое правдоподобие, ABC/сопоставление с историей, суррогат и глобальная
оптимизация, дельфийская калибровка & Совместная байесовская подгонка количественно оценила бы гребни, но
сделала бы заголовочные числа зависящими от априорных распределений. \\
Валидация & Формальной нет; распределение стадий --- единственная информативная внутренняя проверка &
Предварительно зарегистрированные критерии, отложенная популяция, набор эталонных тестов, периодический
ретротест, проверка на уровне решений & Предварительная регистрация процедурна и бесплатна, и она --- самый
крупный разрыв между этой моделью и защитимой. \\
Чувствительность & Одна ось (масштаб времени до диагностики), общее зерно, сообщаются кривые & Перебор по
всем «мягким» параметрам, глобальное разложение дисперсии, вероятностный анализ, сценарии, ценность
информации, структурные замены & Перебор по всем «мягким» параметрам показал бы, что время до диагностики ---
не единственное непроверенное предположение. \\
\end{longtable}
}

\section{Где читать дальше}

Этот документ --- «почему» модели. Два смежных документа охватывают остальные ракурсы:

\begin{itemize}[nosep]
  \item \id{model\_for\_dummies.pdf} --- та же модель простым языком, для читателя, которому важно знать, что
        делает программа, а не почему она так устроена.
  \item \id{python\_port\_analysis.pdf} --- технический и статистический разбор: переменные состояния,
        статистические исправления, сделанные при переносе, производительность и методологические улучшения,
        необходимые для формальной модели уровня принятия политических решений.
\end{itemize}

Файл параметров (\id{config/parameters.toml}) содержит комментарий к каждому заданному вручную параметру с
указанием того, что это и оценивается ли он; разделы выше --- это развёрнутая форма этих комментариев. Два
структурных предположения, которые нужно держать в виду при чтении любого результата, --- принятый масштаб
доклинического времени и принятые вероятности излечения, поскольку именно эти две величины данные не могут
закрепить, и именно их анализ чувствительности в его нынешней одноосевой форме охватывает не полностью.

\end{document}